% !TeX root = ../大学生数学竞赛概率论与数理统计试题选编与专题训练.tex
% 由主文档通过 \input 载入；请编译主文档。
\setcounter{chapter}{3} % 保留原书章号。
\chapter{课程衔接训练：题目}
本章收录264道既有课程训练题，全部与竞赛真题分开计数。基础概率题用于熟练计算，统计题用于补齐竞赛推断题所需的抽样分布、估计和检验知识。


\clearpage\section{事件、计数与几何概率}

\begin{exercise}[事件分解与德摩根律]\label{q:Y1.1}
验证事件的运算公式：
\[
A\cup B=A+\overline{A}B,\quad A=AB+A\overline{B}
\]
\[\begin{aligned}\overline {\bigcup\limits_{1 \leqslant j \leqslant n} {{A_j}} }  &= \bigcap\limits_{1 \leqslant j \leqslant n} {{{\overline A }_j}},\quad  \overline {\bigcap\limits_{1 \leqslant j \leqslant n} {{A_j}} }  \\&= \bigcup\limits_{1 \leqslant j \leqslant n} {{{\overline A }_j}} \end{aligned}\]
\[\begin{aligned}\overline {\bigcup\limits_{j \geqslant 1} {{A_j}} }  &= \bigcap\limits_{j \geqslant 1} {{{\overline A }_j}},\quad  \overline {\bigcap\limits_{j \geqslant 1} {{A_j}} }  \\&= \bigcup\limits_{j \geqslant 1} {{{\overline A }_j}} \end{aligned}\]\par\smallskip {\small 详解见第~\pageref{a:Y1.1} 页。}\end{exercise}

\begin{exercise}[用集合语言表示事件]\label{q:Y1.2}
用集合的运算表达以下事件

（a）$A$发生、$B$和$C$不发生；（b）$A$不发生、$B$和$C$发生；（c）$A,B,C$都不发生；（d）仅有$A,B,C$之一发生；（e）至少$A,B,C$之一发生.\par\smallskip {\small 详解见第~\pageref{a:Y1.2} 页。}\end{exercise}

\begin{exercise}[有放回抽球与样本空间]\label{q:Y1.3}
袋中有大小和质量相同的红、黄、蓝、白球各一个，从中有放回地任取三次，每次一个，并记录下颜色.

（a）写出试验的样本空间;

（b）求取到的三个球颜色互不相同的概率.\par\smallskip {\small 详解见第~\pageref{a:Y1.3} 页。}\end{exercise}

\begin{exercise}[将事件列化为互斥事件列]\label{q:Y1.4}
设$\{A_j:j=1,2,\cdots\}$是事件列，求互不相容的$\{B_j:j=1,2,\cdots\}$使得对于$n\geqslant 1$，$\bigcup\limits_{j = 1}^n {{A_j}}  = \bigcup\limits_{j = 1}^n {{B_j}} $和$\bigcup\limits_{j = 1}^\infty  {{A_j}}  = \bigcup\limits_{j = 1}^\infty  {{B_j}} $成立.\par\smallskip {\small 详解见第~\pageref{a:Y1.4} 页。}\end{exercise}

\begin{exercise}[两阶段抽样的对称性]\label{q:Y1.5}
$100$件产品中有$3$件次品.

（a）从中任取$2$件，求至少得一件次品的概率；

（b）从中任取$10$件，再从这$10$件中任取$2$件，求至少得一件次品的概率.\par\smallskip {\small 详解见第~\pageref{a:Y1.5} 页。}\end{exercise}

\begin{exercise}[无放回试钥匙]\label{q:Y1.7}
钥匙串上的$5$把钥匙中只有$1$把可以开房门，现在无放回地试开房门.计算

（a）第三次打开房门的概率；

（b）三次内打开房门的概率；

（c）如果$5$把钥匙中有$2$把可以打开房门，求三次内打开房门的概率.\par\smallskip {\small 详解见第~\pageref{a:Y1.7} 页。}\end{exercise}

\begin{exercise}[生日相同的概率]\label{q:Y1.8}
设每个人的生日随机落在$365$天中的任一天，求$n$个人的生日互不相同的概率和至少有两个人生日相同的概率. 假定各人的生日相互独立；分别讨论 $0\le n\le365$ 与 $n>365$。\par\smallskip {\small 详解见第~\pageref{a:Y1.8} 页。}\end{exercise}

\begin{exercise}[恰好得到一副手套]\label{q:Y1.9}
有 $6$ 副可区分的手套，每副恰有匹配的左、右手套各一只。从共 $12$ 只手套中等可能地选取 $4$ 只，求恰好含一副匹配手套的概率。\par\smallskip {\small 详解见第~\pageref{a:Y1.9} 页。}\end{exercise}

\begin{exercise}[两次抽样的交集]\label{q:Y1.10}
从标有$1$至$n$的$n$个球中任取$m$个，记下号码后放回.再从这$n$个球中任取$k$个，记下号码.求两组号码中恰有$c$个号码相同的概率,\par\smallskip {\small 详解见第~\pageref{a:Y1.10} 页。}\end{exercise}

\begin{exercise}[单位正方形中的圆域]\label{q:Y1.11}
在区间$[0,1]$中随机地取两点，求它们的平方和小于$1$的概率. 假定两点独立且各服从 $U(0,1)$。\par\smallskip {\small 详解见第~\pageref{a:Y1.11} 页。}\end{exercise}

\begin{exercise}[到达时差与等待]\label{q:Y1.12}
二人相约在$5:00-6:00$点之间见面，如果两人在$5:00-6:00$之间独立地随机到达，求一人至少等另一人半小时的概率.\par\smallskip {\small 详解见第~\pageref{a:Y1.12} 页。}\end{exercise}

\begin{exercise}[硬币落格与边界区域]\label{q:Y1.13}
一枚直径为 $1$ 的硬币落在边长为 $a>0$ 的无限方格纸上。假定圆心在一个基本方格内均匀分布。求硬币不与任何格线相交的概率，并求使该概率小于 $0.01$ 的 $a$ 的范围。\par\smallskip {\small 详解见第~\pageref{a:Y1.13} 页。}\end{exercise}

\begin{exercise}[两次服务是否相遇]\label{q:Y1.14}
在同一个小时内有两辆汽车独立到达同一个加油站加油，车$A$加油需要$5$min，车$B$需要$8$min.如果每辆车在这一个小时内等可能的到达，计算这两辆车在加油站不能相遇的概率.\par\smallskip {\small 详解见第~\pageref{a:Y1.14} 页。}\end{exercise}

\begin{exercise}[立方体内的球域]\label{q:Y1.15}
在$[0,1]$中任取三点$X,Y,Z$，求$X^2+Y^2+Z^2\leqslant 1$的概率和$X^2+Y^2+Z^2=1$的概率. 假定三点相互独立且各服从 $U(0,1)$。\par\smallskip {\small 详解见第~\pageref{a:Y1.15} 页。}\end{exercise}

\begin{exercise}[概率的次可加性]\label{q:Y1.17}
验证概率的次可加性:
\[P(\bigcup\limits_{j = 1}^\infty  {{A_j}) \leqslant } \sum\limits_{j = 1}^\infty  {P({A_j})} \]\par\smallskip {\small 详解见第~\pageref{a:Y1.17} 页。}\end{exercise}

\begin{exercise}[可数个概率一事件]\label{q:Y1.18}
如果事件$A_1,A_2,\cdots$发生的概率都是$1$，证明：
\[P(\bigcap\limits_{j = 1}^\infty  {{A_j}} ) = 1\quad P(\bigcap\limits_{j = 1}^n {{A_j}} ) = 1\]\par\smallskip {\small 详解见第~\pageref{a:Y1.18} 页。}\end{exercise}

\begin{exercise}[四个互斥区域的分割]\label{q:Y1.19}
对于事件$A,B$，验证
\[
P(AB)+P(\overline{A}B)+P(A\overline{B})+P(\overline{A}\:\overline{B})=1
\]\par\smallskip {\small 详解见第~\pageref{a:Y1.19} 页。}\end{exercise}

\begin{exercise}[递增事件列的极限]\label{q:Y1.20}
设$A_1\subset A_2\subset \cdots$，$P(A_j)=(1-1/j)^j$.计算$P(\bigcup\limits_{j = 1}^\infty  {{A_j}} )$.\par\smallskip {\small 详解见第~\pageref{a:Y1.20} 页。}\end{exercise}

\clearpage\section{条件概率、独立性与贝叶斯公式}

\begin{exercise}[至少一次求职成功]\label{q:Y2.1}
甲毕业时在单位A求职成功的概率为$0.4$，在单位B求职成功的概率为$0.6$，在这两家都求职成功的概率为$0.24$.如果甲在这两个单位都求职，计算他至少被一家录用的概率.\par\smallskip {\small 详解见第~\pageref{a:Y2.1} 页。}\end{exercise}

\begin{exercise}[自习室占用问题]\label{q:Y2.2}
$9$个学生随机走进$3$个自习室.

（a）求至少有一个自习室无学生的概率；

（b）求每个自习室至少有一个学生的概率. 假定每位学生独立且等概率选择每个自习室。\par\smallskip {\small 详解见第~\pageref{a:Y2.2} 页。}\end{exercise}

\begin{exercise}[无人拿到自己帽子]\label{q:Y2.3}
$6$个人将各自的帽子混在一起后每人任取一项，求没有人取得自己帽子的概率.\par\smallskip {\small 详解见第~\pageref{a:Y2.3} 页。}\end{exercise}

\begin{exercise}[有放回点名的覆盖概率]\label{q:Y2.4}
开学时老师声明本学期要点名$m$个人，每次都是有放回的随机点名.当全班有$n$个学生，计算

（a）你在本学期不被点到的概率；

（b）当$n\leqslant m$，至少有一个学生在本学期不被点到的概率.\par\smallskip {\small 详解见第~\pageref{a:Y2.4} 页。}\end{exercise}

\begin{exercise}[信封匹配与容斥]\label{q:Y2.5}
将$n$个不同的信笺随机放入$n$个写好地址的信封，求至少有一封匹配的概率.\par\smallskip {\small 详解见第~\pageref{a:Y2.5} 页。}\end{exercise}

\begin{exercise}[独立任务的共同成功]\label{q:Y2.6}
$6$个人独立破译同一个密码.当第$j$个人能破译密码的概率为$p_j$，计算密码被破译的概率.\par\smallskip {\small 详解见第~\pageref{a:Y2.6} 页。}\end{exercise}

\begin{exercise}[发现信号的试验次数]\label{q:Y2.7}
在旷野迷路后，甲每隔一小时发出一个求救信号，如果每个信号被搜救队发现的概率是$0.1$，要以$0.95$的概率保证搜救队能发现信号，甲至少要发送多少个信号？ 假设各次信号是否被发现相互独立。\par\smallskip {\small 详解见第~\pageref{a:Y2.7} 页。}\end{exercise}

\begin{exercise}[保证成功概率的人员数]\label{q:Y2.9}
$12$小时内，每个破译员能够成功破译刚截获的密码的概率（称为破译能力）不小于$0.05$，现在上级派来$5$位破译能力达到$0.2$的专家参加工作，需要再配备多少破译员才能使得密码被破译的概率达到$0.95$. 假设各人的破译结果相互独立。\par\smallskip {\small 详解见第~\pageref{a:Y2.9} 页。}\end{exercise}

\begin{exercise}[德梅雷的骰子问题]\label{q:Y2.10}
投掷两个骰子$24$次，至少掷出一对$6$的概率小于二分之一，证明这一结论 假定骰子公平且所有投掷独立。\par\smallskip {\small 详解见第~\pageref{a:Y2.10} 页。}\end{exercise}

\begin{exercise}[条件成绩概率]\label{q:Y2.11}
无论哪门课先考，甲的微积分和线性代数得优的概率分别是$a$和$b$，微积分先考并得优的条件下，线性代数得优的概率是$c$.当微积分先考时，计算

（a）甲的这两门都得优的概率；（b）微积分没得优的条件下，线性代数得优的概率；（c）线性代数得优的概率；（d）至少一门得优的概率；（e）至少一门不得优的概率. 设 $0<a<1$，且给定概率满足联合分布的相容条件。\par\smallskip {\small 详解见第~\pageref{a:Y2.11} 页。}\end{exercise}

\begin{exercise}[连续受损的生存概率]\label{q:Y2.12}
一架相机第一次落地摔坏的概率是$0.5$.若第一次没摔坏，第二次落地摔坏的概率是$0.7$.若第二次没摔坏，第三次落地摔坏的概率为$0.9$.求该相机落地三次没有摔坏的概率.\par\smallskip {\small 详解见第~\pageref{a:Y2.12} 页。}\end{exercise}

\begin{exercise}[已知一枚骰子的点数]\label{q:Y2.13}
投掷$2$个均匀的骰子，已知至少有一个骰子的点数是$3$，

（a）计算这两个骰子的点数之和为$7$的概率；（b）你应当猜这两个骰子的点数之和是$7$还是$6$.\par\smallskip {\small 详解见第~\pageref{a:Y2.13} 页。}\end{exercise}

\begin{exercise}[两孩家庭中的条件事件]\label{q:Y2.15}
在一个抽象家庭模型中，两名年龄不同孩子的性别相互独立，男孩概率为 $0.51$，女孩概率为 $0.49$。已知至少有一个女孩，求两人都是女孩的概率。\par\smallskip {\small 详解见第~\pageref{a:Y2.15} 页。}\end{exercise}

\begin{exercise}[通知渠道的全概率]\label{q:Y2.16}
老板有一个不很负责的秘书.当要秘书通知张经理$5$小时后见面时，秘书马上办理，但是只用某种方式通知一次.设秘书用传真通知的概率是$0.3$，用短信通知的概率是$0.2$，用电子邮件通知的概率为$0.5$，而张经理在$5$小时内能收到传真的概率是$0.8$,能看到短信的概率是$0.9$，能看到电子邮件的概率是$0.4$.

（a）计算张经理收到通知的概率;

（b）如果收到通知的张经理也有$5$\%的概率不能来见老板，计算老板不能按时见到张经理的概率；\par\smallskip {\small 详解见第~\pageref{a:Y2.16} 页。}\end{exercise}

\begin{exercise}[隐私调查的随机化回答]\label{q:Y2.17}
在调查家庭暴力所占家庭的比例$p$时，为得到实际的$p$同时又不侵犯个人隐私，调查人员将袋中放入比例是$p_0$的红球和比例是$q_0=1-p_0$的白球.被调查者在袋中任取一球窥视后放回，并承诺取得红球就讲真话，取到白球就讲假话.被调查者只需在匿名调查表中选“是”或“否”，然后将表放入投票箱.没人能知道被调查者是否讲真话和回答的是什么.

（a）如果声称有家庭暴力的家庭比例是$p_1$，求$p$；

（b）如果袋中装有$30$个红球，$50$个白球，调查了$320$个家庭，其中有$195$个家庭回答了是，估算$p$； 假设取球与家庭状态独立，且 $p_0\ne1/2$。\par\smallskip {\small 详解见第~\pageref{a:Y2.17} 页。}\end{exercise}

\begin{exercise}[比分差达到二的比赛]\label{q:Y2.19}
甲乙独立地比赛，每局甲胜的概率为 $p\in(0,1)$，无平局。每两局构成一轮；若同一人在这一轮连胜两局则此人获胜，否则重新开始。求甲最终获胜的概率。\par\smallskip {\small 详解见第~\pageref{a:Y2.19} 页。}\end{exercise}

\begin{exercise}[加强抽样的波利亚球袋]\label{q:Y2.20}
袋中起初有 $b$ 个黑球和 $r$ 个红球，$b,r$ 为正整数。每次等可能取一球后放回，再增加同色球一个。记 $B_i$ 为第 $i$ 次取黑球。求 $P(B_1B_2B_3)$、$P(B_2)$、$P(B_3)$，并用归纳法求 $P(B_n)$。\par\smallskip {\small 详解见第~\pageref{a:Y2.20} 页。}\end{exercise}

\begin{exercise}[已知正点反推交通方式]\label{q:Y2.21}
甲乘汽车、火车的概率分别为$0.6,0.4$.汽车和火车正点到达的概率分别是$0.8$和$0.9$.现在甲已经正点到达，求甲乘的是火车的概率.\par\smallskip {\small 详解见第~\pageref{a:Y2.21} 页。}\end{exercise}

\begin{exercise}[观察合格品后的状态]\label{q:Y2.22}
一台机床工作状态良好时，产品的合格率为$99\%$，机床发生故障时的产品合格率为$50\%$.设每次新开机器时机床处于良好状态的概率为$95\%$.如果新开机器后生产的第一件产品是合格品，判断机器处于良好状态的概率.\par\smallskip {\small 详解见第~\pageref{a:Y2.22} 页。}\end{exercise}

\begin{exercise}[低先验概率下的连续检测]\label{q:Y2.23}
一种新方法检测出某种疾病的概率是百分之百，但是把没病的人判定有病的概率达到百分之五.设群体中该病的发病率是万分之一.

（a）甲在身体普查中被判断患病，甲的确患病的概率是多少？

（b）甲再次复查又被判断有病时，甲的确患病的概率是多少？ 本题为抽象模型，假定两次检测结果在给定是否患病的条件下相互独立。\par\smallskip {\small 详解见第~\pageref{a:Y2.23} 页。}\end{exercise}

\begin{exercise}[目击证据与先验比例]\label{q:Y2.25}
某城市夏利牌出租车占$85\%$，富康牌出租车占$15\%$.这两种出租车都是红色，富康出租车略大一些，每辆车肇事的概率相同.在一次出租车肇事逃逸案件中，有人指证富康车肇事.为了确定是否富康车肇事，在肇事地点和相似的能见度下警方对证人辨别出租车的能力进行了测验，发现证人正确识别富康车的概率是$90\%$，正确认识夏利车的概率是$80\%$.

（a）如果证人没有撒谎，求肇事者是富康车的概率；

（b）如果两个情况相同的证人独立地作出相同的指证，求肇事者是富康车的概率. 多名证人的识别结果假设在给定肇事车型后条件独立。\par\smallskip {\small 详解见第~\pageref{a:Y2.25} 页。}\end{exercise}

\begin{exercise}[无穷多次成功的条件]\label{q:Y2.26}
甲不断地向一个目标射击，第$j$次击中目标的概率为$j^{-\alpha}$，其中$\alpha$是正的常数.如果每次能否击中是相互独立的，当$\alpha$取何值时，甲能够击中目标无穷多次.\par\smallskip {\small 详解见第~\pageref{a:Y2.26} 页。}\end{exercise}

\begin{exercise}[最终不再失败的条件]\label{q:Y2.27}
为了得到某种新材料，试验在条件不断改进的情况下进行.如果第$j$次试验成功的概率是$1-j^{-\alpha}$，其中$\alpha$是正常数.计算某次试验后的所有试验不再失败的概率$P$. 假设各次试验是否成功相互独立。\par\smallskip {\small 详解见第~\pageref{a:Y2.27} 页。}\end{exercise}

\begin{exercise}[加红球后的黑球出现次数]\label{q:Y2.28}
设想口袋中只有一个黑球和一个红球.现在在袋中每次任取一个球后放回，且第$j$次取球后再加入红球$j$个.如果一直做下去，能够取到黑球多少次？\par\smallskip {\small 详解见第~\pageref{a:Y2.28} 页。}\end{exercise}

\begin{exercise}[黑红同时增长的抽球]\label{q:Y2.29}
袋中最初有一个黑球和一个红球。每次随机取一球后放回，并在第 $j$ 次抽取后确定地增加 $j$ 个红球和一个黑球。各次使用独立的随机抽取机制。证明黑球几乎必然被抽到无穷多次。\par\smallskip {\small 详解见第~\pageref{a:Y2.29} 页。}\end{exercise}

\begin{exercise}[全概率公式的证明]\label{q:Y2.30}
证明全概率公式：如果事件$A_1,A_2,\cdots$互不相容，$B\subset \bigcup\limits_{j = 1}^\infty  {{A_j}} $，有
\[P(B) = \sum\limits_{j = 1}^\infty  {P({A_j})P(B|{A_j})} \] 对出现在条件概率中的事件，假设 $P(A_j)>0$。\par\smallskip {\small 详解见第~\pageref{a:Y2.30} 页。}\end{exercise}

\begin{exercise}[贝叶斯公式的证明]\label{q:Y2.31}
证明贝叶斯公式：如果事件$A_1,A_2,\cdots$互不相容，$B\subset \bigcup\limits_{j = 1}^\infty  {{A_j}} $，则$P(B)>0$时，有
\[P({A_j}|B) = \frac{{P({A_j})P(B|{A_j})}}{{\sum\limits_{i = 1}^\infty  {P({A_i})P(B|{A_i})} }},j \geqslant 1\] 对出现在条件概率中的事件，假设 $P(A_j)>0$。\par\smallskip {\small 详解见第~\pageref{a:Y2.31} 页。}\end{exercise}

\begin{exercise}[串联与并联可靠性]\label{q:Y2.32}
设元件$A$和$B$在一年内烧断的概率分别是$p_1$和$p_2$.设它们是否烧断是相互独立的.在一年内，

（a）当$A,B$并联时，求$A,B$构成的系统断电的概率；

（b）当$A,B$串联时，求$A,B$构成的系统断电的概率.\par\smallskip {\small 详解见第~\pageref{a:Y2.32} 页。}\end{exercise}

\begin{exercise}[比赛局数与胜率]\label{q:Y2.33}
如果你的水平高于对手，为保证比赛的胜利，你期望三局两胜的比赛规则还是五局三胜的比赛规则？ 设各局结果独立，每局获胜概率均为 $p\in(1/2,1)$，无平局。\par\smallskip {\small 详解见第~\pageref{a:Y2.33} 页。}\end{exercise}

\begin{exercise}[集体迟到的提前量]\label{q:Y2.34}
设开会时每个人是否迟到是相互独立的.如果每个人迟到至少$t$min的概率是$(t+1)^{-2}/100$，在组织一个有$1000$人参加的大会时，要以$0.99$的概率保证无人迟到，应当通知提前多长时间入场.\par\smallskip {\small 详解见第~\pageref{a:Y2.34} 页。}\end{exercise}

\clearpage\section{随机变量及其分布}

\begin{exercise}[分布函数的无穷远极限]\label{q:Y3.1}
用概率的连续性证明分布函数的性质：
\[\begin{aligned}\mathop {\lim }\limits_{x \to \infty } F(x) &= F(\infty ) \\&= 1,\mathop {\lim }\limits_{x \to  - \infty } F(x) \\&= F( - \infty ) \\&= 0\end{aligned}\]\par\smallskip {\small 详解见第~\pageref{a:Y3.1} 页。}\end{exercise}

\begin{exercise}[指数型密度的归一化]\label{q:Y3.2}
设$g(x)=\lambda e^{-3x},x\geqslant 0$是概率密度，求常数$\lambda$.\par\smallskip {\small 详解见第~\pageref{a:Y3.2} 页。}\end{exercise}

\begin{exercise}[二次指数可积的条件]\label{q:Y3.3}
如果常数$a,b,c$使得$g(x)=\exp(ax^2+bx+c),x\in (-\infty,\infty)$是概率密度，则必有$a<0$.\par\smallskip {\small 详解见第~\pageref{a:Y3.3} 页。}\end{exercise}

\begin{exercise}[最大值与最小值恒等式]\label{q:Y3.4}
设$X,Y$是随机变量，验证

（a）$\max(X,Y)=(|X-Y|+X+Y)/2$；

（b）$\min(X,Y)=(X+Y-|X-Y|)/2$；

（c）$\max(X,Y)+\min(X,Y)=X+Y$.\par\smallskip {\small 详解见第~\pageref{a:Y3.4} 页。}\end{exercise}

\begin{exercise}[分段密度的归一化]\label{q:Y3.5}
设$X$有概率密度$f(x) = \left\{ \begin{gathered}
  (x - 1)/a \hfill \quad x\in[1,2]\\
  (x - 3)/a \hfill \quad x\in[3,5]\\
  0 \hfill \quad \text{其他}\\ 
\end{gathered}  \right.$，计算$a$.\par\smallskip {\small 详解见第~\pageref{a:Y3.5} 页。}\end{exercise}

\begin{exercise}[二项计数的尾概率]\label{q:Y3.6}
设某手机一天收到了$8$个短信，每个短信是公事短信的概率是$0.2$，是私事短信的概率是$0.8$.用$X$表示这天收到的公事短信数，计算$X$的概率分布和$P(X\geqslant 3)$. 假设各短信类别相互独立。\par\smallskip {\small 详解见第~\pageref{a:Y3.6} 页。}\end{exercise}

\begin{exercise}[两组二项计数的比较]\label{q:Y3.7}
甲、乙击中目标的概率分别是$0.6,0.7$，如果各射击$3$次，

（a）计算击中次数相同的概率；

（b）计算甲击中的次数多的概率. 假设全部射击结果相互独立。\par\smallskip {\small 详解见第~\pageref{a:Y3.7} 页。}\end{exercise}

\begin{exercise}[首次与第多次成功]\label{q:Y3.8}
编辑正在校对一本书稿.假设该书稿的每页有一个打印错误的概率是$0.05$，没有打印错误的概率是$0.95$.计算编辑

（a）在第$k$页发现第一个打印错误的概率

（b）在第$k$页发现第$r$个打印错误的概率

（c）发现第$r$个错误前，恰有$k$页无打印错误的概率

（d）如果这本书一共有$n$页，计算全书打印错误数的概率分布 假设各页错误情况相互独立，且每页至多一个错误。\par\smallskip {\small 详解见第~\pageref{a:Y3.8} 页。}\end{exercise}

\begin{exercise}[泊松分裂与独立性]\label{q:Y3.9}
甲每天收到的电子邮件数服从泊松分布$P(\lambda)$，且每封电子邮件被过滤掉的概率是$0.2$.

（a）当有$n$封电子邮件发给甲，计算他见到其中$k$封的概率$h_k$；

（b）计算甲每天见到的电子邮件数的分布；

（c）已知甲看到了自己的$k$封电子邮件，计算他有$m$封被过滤掉的概率；

（d）甲每天见到的邮件数和被滤掉的邮件数是否独立. 给定邮件总数后，假设各封邮件独立地被过滤。\par\smallskip {\small 详解见第~\pageref{a:Y3.9} 页。}\end{exercise}

\begin{exercise}[周期信号的均匀等待]\label{q:Y3.13}
求救者每间隔$2$min发出一次瞬时呼叫，随机到达的救援者在收到呼叫信号的范围内至少停留多长时间才能以$0.95$的概率收到呼叫.\par\smallskip {\small 详解见第~\pageref{a:Y3.13} 页。}\end{exercise}

\begin{exercise}[恒定风险率与指数寿命]\label{q:Y3.14}
一个使用了$t$小时的热敏电阻在$\Delta t$内失效的概率是$\lambda \Delta t+o(\Delta t)$，设该热敏电阻的使用寿命是连续性随机变量，求该热敏电阻的寿命分布. 此处的失效概率指已存活到 $t$ 时的条件概率，且初始寿命大于零。\par\smallskip {\small 详解见第~\pageref{a:Y3.14} 页。}\end{exercise}

\begin{exercise}[正态尺寸的合格概率]\label{q:Y3.15}
一台机床加工的部件长度服从正态分布$N(8,36\times 10^{-8})$.当部件的长度在$8\pm 0.0015$内为合格品，求一部件为合格品的概率.\par\smallskip {\small 详解见第~\pageref{a:Y3.15} 页。}\end{exercise}

\begin{exercise}[泊松变量的平方根]\label{q:Y3.16}
设$X\sim P(\lambda)$，计算$Y=\sqrt{X}$的概率分布.\par\smallskip {\small 详解见第~\pageref{a:Y3.16} 页。}\end{exercise}

\begin{exercise}[均匀电流的平方变换]\label{q:Y3.17}
设电流$I$在$8\sim 9$A之间均匀分布.当电流通过$2\Omega$的电阻时，消耗的功率（单位：W）是$W=2I^2$，求$W$的概率密度.\par\smallskip {\small 详解见第~\pageref{a:Y3.17} 页。}\end{exercise}

\begin{exercise}[指数变量的仿射变换]\label{q:Y3.18}
设$X\sim \varepsilon(\lambda)$，$Y=aX+b,a>0$，求$Y$的概率密度.\par\smallskip {\small 详解见第~\pageref{a:Y3.18} 页。}\end{exercise}

\begin{exercise}[半柯西变量的对数]\label{q:Y3.19}
$X$有概率密度
\[f(x)=\frac{2}{\pi(1+x^2)},x\geqslant 0\]
求$Y=\ln X$的概率密度.\par\smallskip {\small 详解见第~\pageref{a:Y3.19} 页。}\end{exercise}

\begin{exercise}[倒数、绝对值与正切变换]\label{q:Y3.20}
设$X$有分段连续的概率密度$f(x)$.

（a）求$Y=X^{-1}$的概率密度；

（b）求$Y=|X|$的概率密度；

（c）求$Y=\tan(X)$的概率密度.\par\smallskip {\small 详解见第~\pageref{a:Y3.20} 页。}\end{exercise}

\begin{exercise}[二项众数与参数估计]\label{q:Y3.21}
设$X$服从二项分布$B(n,p)$.

（a）已知$n=19,p=0.7$，求$p_k=P(X=k)$的最大值点$k$；

（b）已知$n=19,X=9$，求使得$P(X=9)$达到最大的$p$.\par\smallskip {\small 详解见第~\pageref{a:Y3.21} 页。}\end{exercise}

\begin{exercise}[迟到人数的众数和估计]\label{q:Y3.22}
全班有$95$个学生，每个学生上课迟到的概率为$p$.假设每个学生是否迟到是相互独立的.

（a）当$p=0.05$时，有多少个学生迟到的概率最大；

（b）如果有$7$个人迟到，你认为$p$是多少？\par\smallskip {\small 详解见第~\pageref{a:Y3.22} 页。}\end{exercise}

\begin{exercise}[泊松众数与似然]\label{q:Y3.23}
设$X$是服从参数是$\lambda$的泊松分布.

（a）已知$\lambda=23.8$，求$p_k=P(X=k)$的最大值点$k$；

（b）已知$X=21$，求使得$P(X=21)$达到最大的$\lambda$.\par\smallskip {\small 详解见第~\pageref{a:Y3.23} 页。}\end{exercise}

\begin{exercise}[计数结果与参数的区别]\label{q:Y3.24}
假设你每天收到的短信数服从参数是$\lambda$的泊松分布.

（a）已知$\lambda=8.9$时，你今天收到多少个短信的概率最大？

（b）如果你今天收到了$12$个短信，你认为$\lambda$应当是多少？\par\smallskip {\small 详解见第~\pageref{a:Y3.24} 页。}\end{exercise}

\begin{exercise}[骰子样本的最大最小值]\label{q:Y3.25}
将一个骰子投掷$n$次，用$m$表示掷得的最小点数，用$M$表示掷得的最大点数，计算：

（a）$P(m=k),1\leqslant k\leqslant 6$；

（b）$P(M=k),1\leqslant k\leqslant 6$；

（c）$P(m=2,M=5)$.\par\smallskip {\small 详解见第~\pageref{a:Y3.25} 页。}\end{exercise}

\begin{exercise}[圆周投影的反正弦密度]\label{q:Y3.26}
在半径为 $R>0$ 的圆周上按弧长均匀取一点，圆心为原点。求该点横坐标 $X$ 的分布函数及密度。\par\smallskip {\small 详解见第~\pageref{a:Y3.26} 页。}\end{exercise}

\begin{exercise}[正弦变换的两条分支]\label{q:Y3.28}
设$X$有概率密度$f(x)=cx/\pi^2,x\in (0,\pi)$.求$Y=\sin X$的概率密度.\par\smallskip {\small 详解见第~\pageref{a:Y3.28} 页。}\end{exercise}

\begin{exercise}[逆分布函数生成正态数]\label{q:Y3.29}
设$X \sim  U(0,1)$，$\Phi^{-1}(x)$是$\Phi(x)$的反函数，求$Y=\Phi^{-1} (X)$的分布函数. 其中 $\Phi$ 是标准正态分布函数。\par\smallskip {\small 详解见第~\pageref{a:Y3.29} 页。}\end{exercise}

\begin{exercise}[分段单调变换]\label{q:Y3.31}
设$X$服从参数为$\lambda$的指数分布，求
\[Y = \left\{ \begin{gathered}
  X \hfill \quad X\geqslant 1\\
  {X^2} \hfill \quad  X<1\\ 
\end{gathered}  \right.\]

的概率密度.\par\smallskip {\small 详解见第~\pageref{a:Y3.31} 页。}\end{exercise}

\clearpage\section{联合分布与变量变换}

\begin{exercise}[联合概率表与边缘分布]\label{q:Y4.3}
设离散型概率分布$p_{ij}=P(X=i,Y=j)$如下表所示，行指标为 $X$ 的取值，列指标为 $Y$ 的取值。

\begin{center}
    
    \begin{tabular}{|c|c|c|c|c|c|} \hline 
 $p_{ij}$&  $1$&$2$&$3$&$4$&$5$\\ \hline 
 $1$&  $0.06$&$0.05$    &$0.04$&$0.01$&$0.02$\\ \hline 
         $2$&    $0.05$&$0.10$&$0.10$&$0.05$&$0.03$\\ \hline 
         $3$& 
      $0.07$&$0.05$& $0.01$&$0.02$&$0.02$\\ \hline
 $4$& $0.05$& $0.02$& $0.01$&$0.01$&$0.03$\\\hline
 $5$& $0.05$& $0.06$& $0.05$& $0.02$&$0.02$\\\hline\end{tabular}
\end{center}

求$X,Y$的边缘分布；\par\smallskip {\small 详解见第~\pageref{a:Y4.3} 页。}\end{exercise}

\begin{exercise}[同参数二项变量之和]\label{q:Y4.6}
设$X\sim B(n,p),Y\sim B(m,p)$，当$X,Y$独立，计算$X+Y$的概率分布.\par\smallskip {\small 详解见第~\pageref{a:Y4.6} 页。}\end{exercise}

\begin{exercise}[函数依赖的离散向量]\label{q:Y4.8}
如果随机向量$(X,Y)$有如下概率分布：
\[
P(X=i,Y=\frac{1}{i})=c,i=1,2,\cdots,8
\]
确定常数$c$并求$X,Y$的概率分布.\par\smallskip {\small 详解见第~\pageref{a:Y4.8} 页。}\end{exercise}

\begin{exercise}[最小值的尾概率]\label{q:Y4.9}
设随机变量$X,Y$独立同分布，证明
\[
P(a<\min (X,Y)\leqslant b)=P^2(X>a)-P^2(X>b)
\]\par\smallskip {\small 详解见第~\pageref{a:Y4.9} 页。}\end{exercise}

\begin{exercise}[三角支撑上的边缘密度]\label{q:Y4.15}
设随机向量$(X,Y)$有联合密度
\[f(x,y) = \left\{ \begin{gathered}
  {e^{ - x}} \hfill \quad 0<y<x\\
  0 \hfill \quad \text{其他}\\ 
\end{gathered}  \right.\]
求$X,Y$的边缘密度.\par\smallskip {\small 详解见第~\pageref{a:Y4.15} 页。}\end{exercise}

\begin{exercise}[随机家庭规模与男孩数]\label{q:Y4.16}
设$p=1-q\in (0,1)$，$0<\alpha<q/p$.假设一个家庭有$n$个小孩的概率是
\[{p_n} = \left\{ \begin{gathered}
  \alpha {p^n} \hfill \quad n\geqslant 1\\
  1 - \alpha p/q \hfill \quad n=0\\ 
\end{gathered}  \right.\]
如果男孩和女孩的出生是等可能的，求一个家庭有$k$个男孩的概率.\par\smallskip {\small 详解见第~\pageref{a:Y4.16} 页。}\end{exercise}

\begin{exercise}[连续变量的离散随机缩放]\label{q:Y4.22}
设随机变量$X,Y$独立，$X$有概率密度$f(x)$，$Y$有离散型概率分布$P(Y=a_i)=p_i>0,i=1,2,\cdots$，若$a_i$均不为零，证明$Z=XY$有概率密度
\[h(z) = \sum\limits_{i = 1}^\infty  {\frac{{{p_i}}}{{|{a_i}|}}f\left( {\frac{z}{{{a_i}}}} \right)} \]\par\smallskip {\small 详解见第~\pageref{a:Y4.22} 页。}\end{exercise}

\begin{exercise}[博克斯—穆勒正态生成]\label{q:Y4.25}
设随机变量$U$和$V$独立，都服从$(0,1)$上的均匀分布.

（a）计算$r=\sqrt{-2\ln U}$和$\theta=2\pi V$的概率密度.

（b）证明$X=r\cos(\theta),Y=r\sin(\theta)$独立，都服从标准正态分布.\par\smallskip {\small 详解见第~\pageref{a:Y4.25} 页。}\end{exercise}

\begin{exercise}[指数变量的极值与比值]\label{q:Y4.30}
设随机变量$X,Y$独立，$X\sim \mathcal{E}(\lambda),Y\sim \mathcal{E}(\mu)$.分别求$U=\min (X,Y),V=\max(X,Y),W=X/Y$的概率密度.\par\smallskip {\small 详解见第~\pageref{a:Y4.30} 页。}\end{exercise}

\begin{exercise}[平方和与平方比的变换]\label{q:Y4.34}
设 $X,Y$ 独立且均服从率为 $1$ 的指数分布。求 $U=X^2+Y^2,V=X^2/Y^2$ 的联合密度。\par\smallskip {\small 详解见第~\pageref{a:Y4.34} 页。}\end{exercise}

\begin{exercise}[线性变换的密度公式]\label{q:Y4.36}
设 $(X,Y)$ 有联合密度 $f$，可逆矩阵 $A=\begin{pmatrix}a&b\\c&d\end{pmatrix}$，逆矩阵的元素记为 $a_1,b_1,c_1,d_1$。令 $(U,V)^\mathsf T=A(X,Y)^\mathsf T$，证明其联合密度为\[g(u,v)=\frac{f(a_1u+b_1v,c_1u+d_1v)}{|\det A|}.\]\par\smallskip {\small 详解见第~\pageref{a:Y4.36} 页。}\end{exercise}

\begin{exercise}[常数与任意变量独立]\label{q:A4.4}
常数 $a$ 是随机变量, 按定义证明 $a$ 与任何随机变量 $Y$ 独立.\par\smallskip {\small 详解见第~\pageref{a:A4.4} 页。}\end{exercise}

\begin{exercise}[联合分布函数的右连续性]\label{q:A4.5}
证明: 对于固定的 $x$, 联合分布函数 $F(x,y)$ 关于 $y$ 单调不减且右连续.\par\smallskip {\small 详解见第~\pageref{a:A4.5} 页。}\end{exercise}

\begin{exercise}[曲边区域上的联合密度]\label{q:A4.11}
设 $a$ 是常数, $(X,Y)$ 有概率密度
    \[f(x,y)=\begin{cases}
        ax^2y, & x^2<y<1, \\
        0, & \text{其他}.
    \end{cases}\]

    求 $X,Y$ 的边缘密度, 说明 $X,Y$ 不独立.\par\smallskip {\small 详解见第~\pageref{a:A4.11} 页。}\end{exercise}

\begin{exercise}[指数寿命的比较]\label{q:A4.14}
设随机变量 $X,Y$ 相互独立, $X\sim\mathcal{E}(\lambda),Y\sim\mathcal{E}(\mu)$, 计算 $P(X>Y)$.\par\smallskip {\small 详解见第~\pageref{a:A4.14} 页。}\end{exercise}

\begin{exercise}[不独立变量的平方独立]\label{q:A4.17}
设 $(X,Y)$ 在 $(-1,1)^2$ 上的联合密度为 $(1+xy)/4$，其余处为零。证明 $X,Y$ 不独立，但 $X^2,Y^2$ 独立。\par\smallskip {\small 详解见第~\pageref{a:A4.17} 页。}\end{exercise}

\begin{exercise}[曲线下均匀点的投影]\label{q:A4.18}
设 $D$ 是非负连续函数 $g(x)$ 与 $x$ 轴所夹的区域, $D$ 的面积 $m(D)\in(0,\infty)$. 设 $(X,Y)$ 在 $D$ 上均匀分布, 求 $X$ 的概率密度.\par\smallskip {\small 详解见第~\pageref{a:A4.18} 页。}\end{exercise}

\begin{exercise}[离散与连续变量之差]\label{q:A4.19}
设随机变量 $X$ 服从二项分布 $\mathcal{B}(n,p)$, $Y$ 服从指数分布 $\mathcal{E}(\lambda)$. 当 $X,Y$ 独立时, 求 $Z=Y-X$ 的分布函数和概率密度.\par\smallskip {\small 详解见第~\pageref{a:A4.19} 页。}\end{exercise}

\begin{exercise}[求和变换后的边缘积分]\label{q:A4.23}
设随机向量 $(X,Y)$ 有联合密度
    \[f(x,y)=\begin{cases}
        \dfrac{1}{2}(x+y)e^{-(x+y)}, & x,y>0 \\
        0 & \text{其他}.
    \end{cases}\]

    求 $Z=X+Y$ 的概率密度.\par\smallskip {\small 详解见第~\pageref{a:A4.23} 页。}\end{exercise}

\begin{exercise}[独立正态的线性组合]\label{q:A4.27}
设 $X_1,\ldots,X_n$ 相互独立，$X_j\sim N(\mu_j,\sigma_j^2)$。求 $Y=\sum_{j=1}^n a_jX_j+b$ 的分布。\par\smallskip {\small 详解见第~\pageref{a:A4.27} 页。}\end{exercise}

\begin{exercise}[二维正态的条件分布]\label{q:A4.28}
设 $(X_1,X_2)$ 联合正态，均值为 $\mu_1,\mu_2$，标准差 $\sigma_1,\sigma_2>0$，相关系数 $|\rho|<1$。记 $z_i=(x_i-\mu_i)/\sigma_i$、$Q=z_1^2-2\rho z_1z_2+z_2^2$，联合密度为\[f(x_1,x_2)=\frac{e^{-Q/[2(1-\rho^2)]}}{2\pi\sigma_1\sigma_2\sqrt{1-\rho^2}}.\]求 $X_2\mid X_1=x_1$ 的条件分布。\par\smallskip {\small 详解见第~\pageref{a:A4.28} 页。}\end{exercise}

\begin{exercise}[尾分布幂次与竞争概率]\label{q:A4.32}
随机变量 $X,Y$ 相互独立, $X$ 有分布函数 $F_X(x)$ 和概率密度 $f_X(x)=F_X'(x)$. 如果 $P(Y>y)=(P(X>y))^\beta$, 其中 $\beta$ 是正常数, 求 $P(X\geq Y)$.\par\smallskip {\small 详解见第~\pageref{a:A4.32} 页。}\end{exercise}

\clearpage\section{数学期望、方差与协方差}

\begin{exercise}[广告信息之间的等待]\label{q:A5.5}
一部手机收到的短信中有 $p=2\%$ 是广告, 你期望相邻的两次广告短信中有多少个不是广告短信?\par\smallskip {\small 详解见第~\pageref{a:A5.5} 页。}\end{exercise}

\begin{exercise}[标准正态半径的期望]\label{q:A5.8}
设 $X,Y$ 相互独立且均服从标准正态分布。求 $R=\sqrt{X^2+Y^2}$ 的密度及数学期望。\par\smallskip {\small 详解见第~\pageref{a:A5.8} 页。}\end{exercise}

\begin{exercise}[交换对称与比例期望]\label{q:A5.9}
设 $X_1,X_2,\cdots,X_n$ 独立同分布, 且 $P(X_i>0)=1$. 对于 $k\leq n$, 计算
    \[E\dfrac{X_1+X_2+\cdots+X_k}{X_1+X_2+\cdots+X_n}.\]\par\smallskip {\small 详解见第~\pageref{a:A5.9} 页。}\end{exercise}

\begin{exercise}[离散极值的尾和公式]\label{q:A5.12}
设 $X_1,\ldots,X_n$ 为相互独立、同分布的正整数值变量，$p_k=P(X_1=k)$。记 $u_k=P(X_1<k)$、$v_k=1-u_k$。证明\[E\min_iX_i=\sum_{k=1}^\infty v_k^n,\qquad E\max_iX_i=\sum_{k=1}^\infty(1-u_k^n).\]\par\smallskip {\small 详解见第~\pageref{a:A5.12} 页。}\end{exercise}

\begin{exercise}[曲边密度的归一化勘误]\label{q:A5.14}
设 $(X,Y)$ 的联合密度为\[f(x,y)=\frac{3}{2x^3y^2},\qquad x>1,\quad 1<xy<x^2,\]其余处为零。验证归一化并求 $EY$。（此处按原解的积分式校订原题密度的幂次。）\par\smallskip {\small 详解见第~\pageref{a:A5.14} 页。}\end{exercise}

\begin{exercise}[全部指数寿命结束的等待]\label{q:A5.17}
设 $5$ 台计算机独立工作, 每台计算机感染病毒前的时间(分别记作 $X_1,\cdots,X_5$)都服从参数为 $\lambda$ 的指数分布 $\mathcal{E}(\lambda)$. 对 $5$ 台计算机都感染病毒前的时间的数学期望是多少?\par\smallskip {\small 详解见第~\pageref{a:A5.17} 页。}\end{exercise}

\begin{exercise}[样本均值与方差的无偏性]\label{q:A5.20}
设 $n\geq 2$, $X_1,X_2,\cdots,X_n$ 独立同分布,
    \[\begin{aligned}\hat{\mu}&=\dfrac{1}{n}\sum\limits_{j=1}^nX_j,\quad\hat{\sigma}^2\\&=\dfrac{1}{n-1}\sum\limits_{j=1}^n(X_j-\hat{\mu})^2,\end{aligned}\]

    验证 $E\hat{\mu}=EX_1,E\hat{\sigma}^2=\operatorname{Var}(X_1)$. 假设 $E(X_1^2)<\infty$。\par\smallskip {\small 详解见第~\pageref{a:A5.20} 页。}\end{exercise}

\begin{exercise}[最优逆方差加权]\label{q:A5.22}
设独立随机变量 $X_1,\ldots,X_n$ 都有均值 $\mu$，方差分别为 $\sigma_i^2>0$。在 $\sum_iw_i=1$ 的约束下，求线性无偏估计量 $\sum_iw_iX_i$ 的最小方差及最优权重。\par\smallskip {\small 详解见第~\pageref{a:A5.22} 页。}\end{exercise}

\begin{exercise}[非负变量的尾积分]\label{q:A5.25}
设非负随机变量 $X$ 有离散分布 $p_j=P(X=x_j),j\geq 1$. 证明:
    
    (1) $P(X>x)=\sum\limits_{j=1}^\infty p_jI[x<x_j]$;

    (2) $EX=\int_0^\infty P(X>x)\mathrm{d}x$.\par\smallskip {\small 详解见第~\pageref{a:A5.25} 页。}\end{exercise}

\begin{exercise}[和与差的相关系数]\label{q:A5.26}
如果 $\operatorname{Cov}(X,Y)=0$, 验证 $\rho(X+Y,X-Y)=\dfrac{\operatorname{Var}(X)-\operatorname{Var}(Y)}{\operatorname{Var}(X)+\operatorname{Var}(Y)}$. 假设二阶矩有限且 $\operatorname{Var}(X)+\operatorname{Var}(Y)>0$。\par\smallskip {\small 详解见第~\pageref{a:A5.26} 页。}\end{exercise}

\begin{exercise}[正方形密度的协方差]\label{q:A5.28}
设 $(X,Y)$ 有概率密度
    \[f(x,y)=\begin{cases}
        x+y & x\in(0,1),y\in(0,1), \\
        0 & \text{其他}.
    \end{cases}\]

    计算 $\operatorname{Cov}(X,Y)$.\par\smallskip {\small 详解见第~\pageref{a:A5.28} 页。}\end{exercise}

\begin{exercise}[正态尺寸差与装配]\label{q:A5.30}
设活塞的直径 $X$ 的平均值为 $20.00\text{cm}$, 标准差为 $0.02$; 气缸的直径 $Y$ 的平均值为 $20.10\text{cm}$, 标准差为 $0.02$. 设 $X$ 与 $Y$ 独立且都服从 正态分布, 计算活塞能装入气缸的概率.\par\smallskip {\small 详解见第~\pageref{a:A5.30} 页。}\end{exercise}

\begin{exercise}[不相关但不独立的反例]\label{q:A5.32}
设随机变量 $X$ 具有偶对称密度，且 $0<EX^2<\infty$。证明 $X$ 与 $|X|$ 不相关但不独立。\par\smallskip {\small 详解见第~\pageref{a:A5.32} 页。}\end{exercise}

\begin{exercise}[中心正态的高阶矩]\label{q:A5.33}
设 $X\sim N(0,\sigma^2)$, $n\in\mathbb{N}_+$, 验证
    \[EX^n=\begin{cases}
        \sigma^n(n-1)!!, & n=2m, \\
        0 & \text{其他}.
    \end{cases}\]\par\smallskip {\small 详解见第~\pageref{a:A5.33} 页。}\end{exercise}

\begin{exercise}[圆周与圆内均匀模型]\label{q:A5.35}
在半径为 $R>0$ 的圆周上按弧长均匀取一点 $(X,Y)$。求 $EX,EY,\operatorname{Var}(X),\operatorname{Var}(Y)$ 及 $\operatorname{Cov}(X,Y)$，并判断独立性。\par\smallskip {\small 详解见第~\pageref{a:A5.35} 页。}\end{exercise}

\clearpage\section{条件期望与生成函数}

\begin{exercise}[随机组数下的条件二项]\label{q:A6.1}
设 $N_1,N_2,\cdots$ 独立同分布, $N_1\sim\mathcal{B}(8,0.6)$, $M\sim\mathcal{P}(10)$, $M$ 与 $\{N_j\}$ 独立. 对于 $T=\sum\limits_{j=1}^MN_j,k\geq1$, 给出 $T|\{M=k\}$ 的概率分布.\par\smallskip {\small 详解见第~\pageref{a:A6.1} 页。}\end{exercise}

\begin{exercise}[几何个指数等待之和]\label{q:A6.3}
设 $N$ 服从参数为 $p$ 的几何分布, $X_1,X_2,\cdots$ 相互独立, 都服从参数为 $\lambda$ 的指数分布, 当 $N$ 和 $\{X_j\}$ 独立时, 计算 $W=\sum\limits_{j=1}^NX_j$ 的数学期望. 这里几何分布取值 $1,2,\ldots$，质量函数为 $p(1-p)^{k-1}$。\par\smallskip {\small 详解见第~\pageref{a:A6.3} 页。}\end{exercise}

\begin{exercise}[复合泊松等待的均值方差]\label{q:A6.4}
某出租车在一天内遇到的红灯数 $N$ 满足参数为 $\lambda$ 的 泊松 分布. 如果在每个红灯处的等待时间相互独立, 都在 $(0,2)$ 中均匀分布, 计算他一天内用于等候红灯的时间 $X$ 的数学期望和方差. 假设红灯数与各次等待时间独立。\par\smallskip {\small 详解见第~\pageref{a:A6.4} 页。}\end{exercise}

\begin{exercise}[条件指数与条件伽马]\label{q:A6.9}
设 $(X,Y)$ 有联合密度 $f(x,y)=ye^{-y(x+1)},x>0,y>0$. 计算 $X|\{Y=y\},Y|\{X=x\}$ 的概率密度 $f_{X|Y}(x|y),f_{Y|X}(y|x)$.\par\smallskip {\small 详解见第~\pageref{a:A6.9} 页。}\end{exercise}

\begin{exercise}[威布尔寿命的剩余等待]\label{q:A6.14}
假设某设备的使用寿命 $Y$ 服从 威布尔（Weibull）分布, 有概率密度
    \[f(y)=aby^{b-1}e^{-ay^b},\quad y\geq0.\]

    对 $b=1/2$, 计算平均寿命 $E(Y-t|Y>t)$. 其中 $a>0$，所求量是已经存活到 $t$ 后的平均剩余寿命。\par\smallskip {\small 详解见第~\pageref{a:A6.14} 页。}\end{exercise}

\begin{exercise}[泊松观察后的伽马后验]\label{q:A6.15}
设 $\alpha,\beta>0$。随机强度 $Y$ 服从 $\operatorname{Gamma}(\alpha,\beta)$ 分布，其中 $\beta$ 为率参数。给定 $Y=y$ 后，计数 $N$ 的条件分布为 $\operatorname{Pois}(y)$。观察到 $N=n$ 后，求 $Y$ 的后验分布和后验均值。\par\smallskip {\small 详解见第~\pageref{a:A6.15} 页。}\end{exercise}

\begin{exercise}[七个接口的第三次失效]\label{q:A6.17}
一台计算机有 $7$ 个 USB 接口, 至少有 $5$ 个接口正常时, 就认为计算机的 USB 接口能用. 假设开始时每个 USB 接口都正常, 其寿命 $X_1,X_2,\cdots,X_7$ 相互独立, 都服从参数为 $\lambda$ 的指数分布. 计算这台计算机的 USB 能用的时间 $Y$ 的概率分布.\par\smallskip {\small 详解见第~\pageref{a:A6.17} 页。}\end{exercise}

\begin{exercise}[指数顺序统计量的间隔]\label{q:A6.20}
设 $X_1,X_2,\cdots,X_n$ 独立同分布, 都服从参数为 $\lambda$ 的指数分布, 证明随机变量
    \[Y_1=X_{(1)},\quad Y_k=X_{(k)}-X_{(k-1)},\quad k=2,3,\cdots,n\]
    相互独立, 且 $Y_i\sim\mathcal{E}((n+1-i)\lambda)$.\par\smallskip {\small 详解见第~\pageref{a:A6.20} 页。}\end{exercise}

\begin{exercise}[仿射变换的概率母函数]\label{q:A6.21}
设取非负整数值的随机变量 $X$ 有母函数 $g(s)$, 对非负整数 $a,b$, 求 $Y=aX+b$ 的母函数 $h(s)$.\par\smallskip {\small 详解见第~\pageref{a:A6.21} 页。}\end{exercise}

\begin{exercise}[两组硬币成功数之差]\label{q:A6.24}
两人各抛均匀的硬币 $n$ 次, 计算甲的正面次数恰好大于乙的正面次数 $k$ 次的概率.\par\smallskip {\small 详解见第~\pageref{a:A6.24} 页。}\end{exercise}

\begin{exercise}[均匀和的特征函数极限]\label{q:A6.29}
设 $X_1,X_2,\cdots$ 独立同分布, 都在 $(-1,1)$ 中均匀分布. 定义
    \[Y_n=\sqrt{\dfrac{3}{n}}\sum\limits_{j=1}^nX_j.\]

    (1) 计算 $X_1$ 的特征函数 $\phi_X(t)$;

    (2) 计算 $Y_n$ 的特征函数 $\phi_{Y_n}(t)$, 证明 $Y_n\xrightarrow{d}N(0,1)$.\par\smallskip {\small 详解见第~\pageref{a:A6.29} 页。}\end{exercise}

\begin{exercise}[指数和的特征函数极限]\label{q:A6.30}
设 $X_1,X_2,\cdots$ 独立同分布, $X_1\sim\mathcal{E}(\lambda)$. 定义
    \[S_n=\sum\limits_{j=1}^nX_j,\quad Y_n=\dfrac{S_n-n/\lambda}{\sqrt{n/\lambda^2}}.\]

    (1) 计算 $X_1$ 的特征函数 $\phi_X(t)$;

    (2) 计算 $Y_n$ 的特征函数 $\phi_{Y_n}(t)$;

    (3) 证明 $Y_n\xrightarrow{d}N(0,1)$.\par\smallskip {\small 详解见第~\pageref{a:A6.30} 页。}\end{exercise}

\begin{exercise}[二项分布的泊松极限]\label{q:A6.34}
设 $X_n\sim\operatorname{Bin}(n,p_n)$，且 $np_n\to\lambda\in[0,\infty)$。证明 $X_n\xrightarrow{d}\operatorname{Pois}(\lambda)$。\par\smallskip {\small 详解见第~\pageref{a:A6.34} 页。}\end{exercise}

\begin{exercise}[给定总和后的部分和]\label{q:A6.35}
设 $X_1,\ldots,X_n$ 独立同分布且 $E|X_1|<\infty$，$S_k=\sum_{j=1}^kX_j$。证明 $E(S_k\mid S_n)=kS_n/n$，$1\le k\le n$。\par\smallskip {\small 详解见第~\pageref{a:A6.35} 页。}\end{exercise}

\clearpage\section{大数定律与中心极限定理}

\begin{exercise}[截尾期望的概率上界]\label{q:A7.2}
证明: 设非负随机变量 $X$ 有概率密度 $f(x)$, 则 $\forall$ 正数 $M$, 有
    \[P(X\geq M)\leq\dfrac{1}{M}\int_M^\infty xf(x)\mathrm{d}x.\]\par\smallskip {\small 详解见第~\pageref{a:A7.2} 页。}\end{exercise}

\begin{exercise}[柯西不等式与尾概率]\label{q:A7.3}
证明: 设非负随机变量 $X$ 有分布函数 $F(x)$, 则 $\forall$ 正数 $M$, 有
    \[P(X>M)\leq\dfrac{\sqrt{(1-F(M))EX^2}}{M}.\] 假设有关二阶矩有限，且 $M>0$。\par\smallskip {\small 详解见第~\pageref{a:A7.3} 页。}\end{exercise}

\begin{exercise}[最小值的概率界]\label{q:A7.4}
设 $X_1,\ldots,X_n$ 独立同分布且非负，$EX_1^2<\infty$。对 $M>0$，证明\[P(\min_iX_i\ge M)\le\left(\frac{EX_1^2\,P(X_1\ge M)}{M^2}\right)^{n/2}.\]\par\smallskip {\small 详解见第~\pageref{a:A7.4} 页。}\end{exercise}

\begin{exercise}[伽马矩与切比雪夫界]\label{q:A7.5}
设 $X$ 有概率密度 $f(x)=\dfrac{x^m}{m!}e^{-x},\ x\geq0$, 计算 $EX,\operatorname{Var}(X)$, 证明
    \[P(0<X<2(m+1))\geq\dfrac{m}{m+1}.\]\par\smallskip {\small 详解见第~\pageref{a:A7.5} 页。}\end{exercise}

\begin{exercise}[三角函数样本均值之比]\label{q:A7.10}
设 $X_1,X_2,\cdots$ 独立同分布, 都在 $(0,\pi/2)$ 中均匀分布,
    \[Y_n=\dfrac{\sin X_1+\sin X_2+\cdots+\sin X_n}{\cos X_1+\cos X_2+\cdots+\cos X_n},\quad n\geq1,\]
    
    当 $n\to\infty$ 时, 计算 $Y_n$ 的 a.s. 极限.\par\smallskip {\small 详解见第~\pageref{a:A7.10} 页。}\end{exercise}

\begin{exercise}[比例调查的样本量]\label{q:A7.14}
利用独立伯努利抽样估计未知比例 $p$。希望样本比例与 $p$ 的误差不超过 $0.01$，置信水平约为 $99\%$。用中心极限定理求不依赖 $p$ 的保守样本量，并说明该结果是否是有限样本的严格保证。\par\smallskip {\small 详解见第~\pageref{a:A7.14} 页。}\end{exercise}

\begin{exercise}[两组样本均值差的近似]\label{q:A7.15}
两组样本相互独立；第一组为来自均值 $\mu_1$、有限正方差 $\sigma_1^2$ 的独立同分布样本，容量 $n$；第二组相应为 $\mu_2,\sigma_2^2,m$。当 $n,m$ 都较大时，求两个样本均值之差的近似分布。\par\smallskip {\small 详解见第~\pageref{a:A7.15} 页。}\end{exercise}

\begin{exercise}[泊松概率和的二分极限]\label{q:A7.20}
证明\[\begin{aligned}\lim_{n\to\infty}e^{-n}\sum_{k=0}^{n}\frac{n^k}{k!}&=\frac12,\qquad\lim_{n\to\infty}e^{-n}\sum_{k=0}^{n-1}\frac{n^k}{k!}\\&=\frac12.\end{aligned}\]\par\smallskip {\small 详解见第~\pageref{a:A7.20} 页。}\end{exercise}

\begin{exercise}[缩放收敛与紧性]\label{q:A7.21}
如果 $\sqrt{n}\xi_n\xrightarrow{d}N(0,1)$, 证明 $\xi_n\xrightarrow{p}0$.\par\smallskip {\small 详解见第~\pageref{a:A7.21} 页。}\end{exercise}

\begin{exercise}[不同参数二项和的极限]\label{q:A7.22}
设 $X_j$ 相互独立且 $X_j\sim\operatorname{Bin}(j,p_j)$，其中 $0.1\le p_j\le0.9$。令 $S_n=\sum_{j=1}^nX_j$。证明 $(S_n-ES_n)/\sqrt{\operatorname{Var}(S_n)}$ 依分布收敛于标准正态分布。\par\smallskip {\small 详解见第~\pageref{a:A7.22} 页。}\end{exercise}

\begin{exercise}[小扰动不改变分布极限]\label{q:A7.25}
如果 $\xi_n\xrightarrow{d}N(0,1),\eta_n\xrightarrow{p}0$, 证明 $\xi_n+\eta_n\xrightarrow{d}N(0,1)$.\par\smallskip {\small 详解见第~\pageref{a:A7.25} 页。}\end{exercise}

\clearpage\section{样本与抽样分布}

\begin{exercise}[统计量与未知参数]\label{q:M-F-C4}
地铁延误时间 $X_1,\ldots,X_5$ 独立同分布，均服从未知率参数 $\lambda>0$ 的指数分布。判断以下量是否为统计量，并说明理由：$1/\lambda$、样本均值 $\overline X$、$\overline X-1/\lambda$、样本标准差 $S$。\par\smallskip {\small 详解见第~\pageref{a:M-F-C4} 页。}\end{exercise}

\begin{exercise}[样本均值的均值与方差]\label{q:B8.1}
设 $X_1,\ldots,X_n$ 独立同分布，$EX_i=\mu$，$\operatorname{Var}(X_i)=\sigma^2<\infty$。求 $E\overline X$、$\operatorname{Var}(\overline X)$ 及 $\operatorname{Cov}(X_i,\overline X)$，并说明若样本相关，哪些公式仍成立。\par\smallskip {\small 详解见第~\pageref{a:B8.1} 页。}\end{exercise}

\begin{exercise}[样本方差的无偏性]\label{q:B8.2}
在独立同分布且方差有限的条件下，令 $S^2=(n-1)^{-1}\sum_{i=1}^n(X_i-\overline X)^2$，$n\ge2$。证明 $ES^2=\sigma^2$，并求使用分母 $n$ 的样本二阶中心矩的偏差。\par\smallskip {\small 详解见第~\pageref{a:B8.2} 页。}\end{exercise}

\begin{exercise}[正态样本均值和残差的独立性]\label{q:B8.3}
设 $X_i$ 独立同分布于 $N(\mu,\sigma^2)$，$\sigma>0$。证明 $\overline X$ 与整个残差向量 $(X_1-\overline X,\ldots,X_n-\overline X)$ 独立，并推出 $\overline X$ 与 $S^2$ 独立。\par\smallskip {\small 详解见第~\pageref{a:B8.3} 页。}\end{exercise}

\begin{exercise}[卡方分布及自由度]\label{q:B8.4}
设 $X_i$ 为容量 $n\ge2$ 的正态样本。证明 $(n-1)S^2/\sigma^2\sim\chi^2_{n-1}$，并解释为何不是 $\chi^2_n$。\par\smallskip {\small 详解见第~\pageref{a:B8.4} 页。}\end{exercise}

\begin{exercise}[学生 t 统计量]\label{q:B8.5}
设 $X_1,\ldots,X_n$ 为来自 $N(\mu,\sigma^2)$ 的样本。推导 $T=\sqrt n(\overline X-\mu)/S$ 的分布，并说明为什么直接将 $S$ 代替 $\sigma$ 后仍使用标准正态分布不是精确做法。\par\smallskip {\small 详解见第~\pageref{a:B8.5} 页。}\end{exercise}

\begin{exercise}[方差比与 F 分布]\label{q:B8.6}
两组相互独立的正态样本容量分别为 $n,m\ge2$，方差为 $\sigma_1^2,\sigma_2^2$。求 $(S_1^2/\sigma_1^2)/(S_2^2/\sigma_2^2)$ 的分布，并推导 $F$ 分位数的倒数关系。\par\smallskip {\small 详解见第~\pageref{a:B8.6} 页。}\end{exercise}

\begin{exercise}[顺序统计量的分布]\label{q:B8.7}
设 $X_i$ 独立同分布于连续分布函数 $F$，密度为 $f$。推导最大值 $X_{(n)}$、最小值 $X_{(1)}$ 的分布函数，以及第 $k$ 个顺序统计量的密度。\par\smallskip {\small 详解见第~\pageref{a:B8.7} 页。}\end{exercise}

\begin{exercise}[均匀样本最大值]\label{q:B8.8}
设 $X_i$ 独立同分布于 $U(0,\theta)$，$\theta>0$。令 $M=X_{(n)}$。求 $EM,\operatorname{Var}(M)$，并构造 $\theta$ 的无偏估计量。\par\smallskip {\small 详解见第~\pageref{a:B8.8} 页。}\end{exercise}

\begin{exercise}[指数样本的最小值与和]\label{q:B8.9}
设 $X_i$ 独立同分布于率为 $\lambda$ 的指数分布。求最小值 $M$ 和总和 $T$ 的分布，并比较 $nM$ 与 $\overline X$ 的均值和方差。\par\smallskip {\small 详解见第~\pageref{a:B8.9} 页。}\end{exercise}

\begin{exercise}[中心极限定理的正确标准化]\label{q:M-F-C13}
设独立同分布样本的均值为 $\mu$，方差为 $0<\sigma^2<\infty$。判断：大样本时每个观测近似正态；样本均值近似正态；$(\overline X-\mu)/\sigma$ 近似标准正态；$\sqrt n(\overline X-\mu)/\sigma$ 近似标准正态。\par\smallskip {\small 详解见第~\pageref{a:M-F-C13} 页。}\end{exercise}

\clearpage\section{矩估计、偏差与一致性}

\begin{exercise}[伯努利矩估计]\label{q:B9.1}
独立样本来自 $\operatorname{Bern}(p)$。用矩法估计 $p$，求估计量的偏差、方差和均方误差，并证明一致性。\par\smallskip {\small 详解见第~\pageref{a:B9.1} 页。}\end{exercise}

\begin{exercise}[均匀分布的矩估计与极值估计]\label{q:B9.2}
样本来自 $U(0,\theta)$。求矩估计 $\widehat\theta_M$，并比较它与 $\widetilde\theta=(n+1)X_{(n)}/n$ 的均方误差。\par\smallskip {\small 详解见第~\pageref{a:B9.2} 页。}\end{exercise}

\begin{exercise}[伽马分布的两个矩参数]\label{q:B9.3}
设样本来自 $\operatorname{Gamma}(\alpha,\beta)$，$\alpha,\beta>0$，$\beta$ 为率。设 $m_1=\overline X$，$v_n=n^{-1}\sum_i(X_i-\overline X)^2$。由前两个矩构造 $\alpha,\beta$ 的估计。\par\smallskip {\small 详解见第~\pageref{a:B9.3} 页。}\end{exercise}

\begin{exercise}[正态分布的矩估计]\label{q:B9.4}
样本来自 $N(\mu,\sigma^2)$，两个参数均未知。用一阶和二阶原点矩求矩估计，并比较其方差估计与无偏样本方差。\par\smallskip {\small 详解见第~\pageref{a:B9.4} 页。}\end{exercise}

\begin{exercise}[偏差与均方误差分解]\label{q:B9.5}
对二阶矩有限的估计量 $T$，证明 $E_\theta(T-\theta)^2=\operatorname{Var}_\theta(T)+[E_\theta T-\theta]^2$。若 $X_i\sim N(\mu,\sigma^2)$ 且 $\sigma^2$ 已知，比较 $\overline X$ 与 $a\overline X$ 的均方误差。\par\smallskip {\small 详解见第~\pageref{a:B9.5} 页。}\end{exercise}

\begin{exercise}[样本方差的最小均方误差缩放]\label{q:B9.6}
正态样本容量 $n\ge2$。在形如 $aS^2$ 的估计量中，求估计 $\sigma^2$ 的最小均方误差系数 $a$。\par\smallskip {\small 详解见第~\pageref{a:B9.6} 页。}\end{exercise}

\begin{exercise}[一致性与渐近正态性]\label{q:B9.7}
设 $X_i$ 独立同分布于 $\operatorname{Pois}(\lambda)$，$\lambda>0$。证明 $\overline X$ 一致估计 $\lambda$，并求 $\log\overline X$ 的渐近分布；解释全零样本的处理。\par\smallskip {\small 详解见第~\pageref{a:B9.7} 页。}\end{exercise}

\begin{exercise}[样本分位数的直观比较]\label{q:B9.8}
对均匀总体 $U(0,\theta)$，比较由 $2\overline X$ 与 $2X_1$ 估计 $\theta$。两者是否无偏、是否一致？由此说明“大样本可用更多数据”与“估计公式实际利用更多数据”的区别。\par\smallskip {\small 详解见第~\pageref{a:B9.8} 页。}\end{exercise}

\clearpage\section{极大似然估计}

\begin{exercise}[离散参数空间中的硬币似然]\label{q:M-C10-1}
独立抛硬币 $80$ 次，观察到 $49$ 次正面、$31$ 次反面。先假定正面概率只可能是 $1/3,1/2,2/3$，求极大似然估计；再允许 $p\in[0,1]$，重新求估计。\par\smallskip {\small 详解见第~\pageref{a:M-C10-1} 页。}\end{exercise}

\begin{exercise}[指数寿命的极大似然估计]\label{q:M-C10-2}
五个独立寿命观测为 $2,3,1,3,4$，假定来自率参数未知的指数分布。写出似然、对数似然和极大似然估计，并给出容量为 $n$ 的一般公式。\par\smallskip {\small 详解见第~\pageref{a:M-C10-2} 页。}\end{exercise}

\begin{exercise}[已知形状的伽马率估计]\label{q:M-E2-II1}
总体密度为 $f(x;\beta)=\beta^3x^2e^{-\beta x}/2$，$x>0$，$\beta>0$。五个独立观测为 $1,1,2,3,5$。求 $\beta$ 的极大似然估计，并推广到已知形状 $\alpha$ 的伽马模型。\par\smallskip {\small 详解见第~\pageref{a:M-E2-II1} 页。}\end{exercise}

\begin{exercise}[正态模型的联合极大似然]\label{q:B10.1}
对正态样本，同时估计未知的 $\mu,\sigma^2$。推导联合极大似然估计，并说明样本所有值相同的特殊情形。\par\smallskip {\small 详解见第~\pageref{a:B10.1} 页。}\end{exercise}

\begin{exercise}[均匀端点与支撑集]\label{q:B10.2}
设样本来自闭区间密度表示 $f(x;\theta)=\theta^{-1}\mathbf1_{[0,\theta]}(x)$。求 $\theta$ 的极大似然估计，并解释只对 $-n\log\theta$ 求导为什么不够。\par\smallskip {\small 详解见第~\pageref{a:B10.2} 页。}\end{exercise}

\begin{exercise}[几何分布的估计与边界]\label{q:B10.3}
设 $P(X=k)=p(1-p)^{k-1}$，$k=1,2,\ldots$，$0<p\le1$。独立样本容量为 $n$，求 $p$ 的极大似然估计。若改成记录首次成功前的失败次数，公式如何改变？\par\smallskip {\small 详解见第~\pageref{a:B10.3} 页。}\end{exercise}

\begin{exercise}[带删失的指数寿命]\label{q:B10.4}
独立寿命 $X_i\sim\operatorname{Exp}(\lambda)$，每个个体最多观察到固定时间 $c>0$。记录 $Y_i=\min(X_i,c)$ 和失效指示 $\delta_i=\mathbf1_{\{X_i\le c\}}$。推导 $\lambda$ 的极大似然估计。\par\smallskip {\small 详解见第~\pageref{a:B10.4} 页。}\end{exercise}

\begin{exercise}[极大似然的不变性]\label{q:B10.5}
若 $\widehat\theta$ 是参数 $\theta$ 的极大似然估计，且 $g$ 为一一映射，证明 $g(\widehat\theta)$ 是 $g(\theta)$ 的极大似然估计。以指数模型的中位数为例应用该结论。\par\smallskip {\small 详解见第~\pageref{a:B10.5} 页。}\end{exercise}

\begin{exercise}[指数率估计的有限样本偏差]\label{q:B10.6}
指数率为 $\lambda$ 的样本总和 $T\sim\operatorname{Gamma}(n,\lambda)$。求 $E(n/T)$，并构造 $\lambda$ 的无偏估计。说明所需样本容量条件。\par\smallskip {\small 详解见第~\pageref{a:B10.6} 页。}\end{exercise}

\clearpage\section{充分统计量与无偏估计}

\begin{exercise}[伯努利样本的充分统计量]\label{q:B11.1}
设 $X_i\sim\operatorname{Bern}(p)$ 独立。证明总成功次数 $T=\sum_iX_i$ 是 $p$ 的充分统计量，并给出给定 $T=t$ 时样本向量的条件分布。\par\smallskip {\small 详解见第~\pageref{a:B11.1} 页。}\end{exercise}

\begin{exercise}[泊松样本的充分性与条件二项分布]\label{q:B11.2}
设 $X_i\sim\operatorname{Pois}(\lambda)$ 独立。证明 $T=\sum_iX_i$ 充分，并求 $X_1\mid T=t$ 的分布。\par\smallskip {\small 详解见第~\pageref{a:B11.2} 页。}\end{exercise}

\begin{exercise}[拉奥—布莱克改进]\label{q:B11.3}
在泊松样本模型下，$U=\mathbf1_{\{X_1=0\}}$ 无偏估计 $e^{-\lambda}$。计算 $E(U\mid T)$，其中 $T=\sum_iX_i$，并证明改进后的均方误差不大于原估计。\par\smallskip {\small 详解见第~\pageref{a:B11.3} 页。}\end{exercise}

\begin{exercise}[二项族的完全性]\label{q:B11.4}
设 $T\sim\operatorname{Bin}(n,p)$，$p\in(0,1)$。若函数 $g$ 满足对所有 $p$ 都有 $E_pg(T)=0$，证明 $g(t)=0$，$t=0,\ldots,n$。说明该结论对最小方差无偏估计的意义。\par\smallskip {\small 详解见第~\pageref{a:B11.4} 页。}\end{exercise}

\begin{exercise}[估计平方概率]\label{q:B11.5}
伯努利样本容量 $n\ge2$。构造 $p^2$ 的无偏估计，并与直接代入估计 $\widehat p^2$ 比较偏差。\par\smallskip {\small 详解见第~\pageref{a:B11.5} 页。}\end{exercise}

\begin{exercise}[伯努利模型的信息量与下界]\label{q:B11.6}
对 $p\in(0,1)$ 的伯努利样本，计算费希尔信息量，给出无偏估计 $p$ 的克拉美—拉奥下界，并判断样本比例是否达到该下界。\par\smallskip {\small 详解见第~\pageref{a:B11.6} 页。}\end{exercise}

\begin{exercise}[信息下界并非总能达到]\label{q:B11.7}
在指数率模型中计算 $I_n(\lambda)$。将克拉美—拉奥下界与 $n>2$ 时无偏估计 $(n-1)/T$ 的方差比较，解释为何二者不必相等。\par\smallskip {\small 详解见第~\pageref{a:B11.7} 页。}\end{exercise}

\begin{exercise}[非正则模型中的错误下界]\label{q:B11.8}
有人对 $U(0,\theta)$ 模型计算 $\partial\log f/\partial\theta=-1/\theta$，据此声称任何无偏估计 $\theta$ 的方差至少为 $\theta^2/n$。指出错误，并给出反证性比较。\par\smallskip {\small 详解见第~\pageref{a:B11.8} 页。}\end{exercise}

\clearpage\section{贝叶斯估计与预测}

\begin{exercise}[三类别观察的贝叶斯更新]\label{q:M-E2-II2}
独立试验产生 $A,B,C$ 三种结果，概率分别为 $\theta/2,\theta/2,1-\theta$，$0<\theta<1$。先验密度为 $3\theta^2$。观察序列 $A,B,C,A,B$ 后，求后验密度及后验均值。\par\smallskip {\small 详解见第~\pageref{a:M-E2-II2} 页。}\end{exercise}

\begin{exercise}[先验赔率与后验赔率]\label{q:M-E2-II3}
某对象具有特征 $H$ 的先验赔率为 $1:100$。筛选规则在 $H$ 下必为阳性，在 $H^c$ 下有 $10\%$ 概率为阳性。阳性后，求后验赔率及后验概率。\par\smallskip {\small 详解见第~\pageref{a:M-E2-II3} 页。}\end{exercise}

\begin{exercise}[贝塔与二项共轭]\label{q:M-F-C14}
成功概率 $\theta$ 的先验为 $\operatorname{Beta}(5,5)$。独立进行 $30$ 次伯努利试验，观察到 $20$ 次成功。求后验分布、后验均值、后验众数，以及下一次成功的预测概率。\par\smallskip {\small 详解见第~\pageref{a:M-F-C14} 页。}\end{exercise}

\begin{exercise}[平方损失下的贝叶斯估计]\label{q:B12.1}
给定数据后，参数 $\Theta$ 的后验二阶矩有限。在平方损失 $(a-\Theta)^2$ 下，证明最优决策为后验均值；并说明绝对损失对应的解。\par\smallskip {\small 详解见第~\pageref{a:B12.1} 页。}\end{exercise}

\begin{exercise}[伽马与泊松共轭的多次观察]\label{q:B12.2}
独立计数 $X_i\mid\lambda\sim\operatorname{Pois}(t_i\lambda)$，其中暴露时间 $t_i>0$ 已知。先验 $\lambda\sim\operatorname{Gamma}(a,b)$，采用率参数。求后验及其均值，并解释暴露时间的作用。\par\smallskip {\small 详解见第~\pageref{a:B12.2} 页。}\end{exercise}

\begin{exercise}[正态均值的共轭更新]\label{q:B12.3}
设 $X_i\mid\mu\sim N(\mu,\sigma^2)$ 独立，$\sigma^2$ 已知；先验 $\mu\sim N(m_0,\tau_0^2)$。求后验分布及下一次观测的预测分布。\par\smallskip {\small 详解见第~\pageref{a:B12.3} 页。}\end{exercise}

\begin{exercise}[主持人规则改变三门问题]\label{q:M-P7-1}
奖品等概率位于 $A,B,C$ 三门之一，选手先选 $A$。观察到主持人打开 $B$ 且里面无奖品。分别求以下规则下的后验概率：主持人知道奖品位置，总打开未选的空门，有两扇可开时等概率选择；主持人不知道位置，随机打开 $B,C$；主持人以 $0.7$ 概率按第一规则、以 $0.3$ 概率按第二规则，且规则选择独立于奖品位置。\par\smallskip {\small 详解见第~\pageref{a:M-P7-1} 页。}\end{exercise}

\begin{exercise}[未知面数骰子的预测]\label{q:M-P7-2}
从面数为 $4,6,8,12,20$ 的五枚公平骰子中等概率选一枚；每枚各面标为 $1$ 至其面数。选定后独立投掷，连续 $n\ge1$ 次都得到 $7$。求面数的后验分布、下一次仍为 $7$ 的预测概率及其极限。\par\smallskip {\small 详解见第~\pageref{a:M-P7-2} 页。}\end{exercise}

\begin{exercise}[显著性与后验概率的区别]\label{q:M-P8-1fg}
独立试验 $250$ 次，正面 $140$ 次。先只考虑 $p=0.5$ 和 $p=0.55$ 两个假设，先验各为 $1/2$，求后者后验概率；再取 $p\sim\operatorname{Beta}(1,1)$，写出后验并求 $P(p>0.5\mid D)$ 的表达式。为什么这些量不是双侧检验的 $p$ 值？\par\smallskip {\small 详解见第~\pageref{a:M-P8-1fg} 页。}\end{exercise}

\clearpage\section{区间估计与样本量}

\begin{exercise}[小样本正态总体的三类区间]\label{q:M-C22-1}
独立正态样本观测为 $4,1,2,3$。先假定总体标准差已知为 $2$，求均值的 $90\%$ 置信区间；再假定标准差未知，求均值、方差和标准差各自的 $90\%$ 置信区间。\par\smallskip {\small 详解见第~\pageref{a:M-C22-1} 页。}\end{exercise}

\begin{exercise}[已知标准差的置信水平比较]\label{q:M-P10-1}
样本容量 $30$，总体为 $N(\theta,9^2)$，样本均值为 $15$。求 $95\%$ 和 $70\%$ 均值置信区间。如果样本容量增至 $90$ 而其他条件不变，区间宽度如何变化？\par\smallskip {\small 详解见第~\pageref{a:M-P10-1} 页。}\end{exercise}

\begin{exercise}[灌装量的正态均值区间]\label{q:M-P10-2}
九个灌装量观测为 $352,351,361,353,352,358,360,358,359$。假设观测独立且总体正态。分别在总体标准差已知为 $3$、总体标准差未知的情况下，求总体均值的 $95\%$ 与 $98\%$ 置信区间。\par\smallskip {\small 详解见第~\pageref{a:M-P10-2} 页。}\end{exercise}

\begin{exercise}[正态方差的区间]\label{q:M-P10-3}
十二个观测为 $6,6.4,7,5.8,6,5.8$；$5.9,6.7,6.1,6.5,6.3,5.8$。假定总体正态，求方差的 $95\%$ 置信区间，并解释区间为何不关于样本方差对称。\par\smallskip {\small 详解见第~\pageref{a:M-P10-3} 页。}\end{exercise}

\begin{exercise}[身高区间与宽度要求]\label{q:M-P11-1}
正态总体的 $20$ 个独立身高观测，样本均值 $69.55$ 英寸，样本方差 $14.26$ 平方英寸。求未知方差时的 $90\%$ 均值区间。若另有可靠信息表明总体标准差为 $3.77$ 英寸，求相应区间，并求使区间总宽度不超过 $1$ 英寸所需样本量。\par\smallskip {\small 详解见第~\pageref{a:M-P11-1} 页。}\end{exercise}

\begin{exercise}[顺序统计量构造端点区间]\label{q:M-P11-2}
从 $U(0,a)$ 中独立抽样 $n\ge2$ 次，记第二小的观测为 $Y_{(2)}$。已知 $Y_{(2)}/a\sim\operatorname{Beta}(2,n-1)$。利用该枢轴构造 $a$ 的 $95\%$ 置信区间；当 $n=9,y_{(2)}=1.5$ 时给出数值。\par\smallskip {\small 详解见第~\pageref{a:M-P11-2} 页。}\end{exercise}

\begin{exercise}[比例误差与样本量]\label{q:M-C22-3}
利用独立伯努利样本估计未知比例 $p$，要求正态近似的 $95\%$ 误差半径不超过 $0.01$。用 $p(1-p)\le1/4$ 规划样本量；改为 $80\%$ 时如何变化？若样本容量固定为 $900$，给出两种置信水平的保守误差半径。\par\smallskip {\small 详解见第~\pageref{a:M-C22-3} 页。}\end{exercise}

\begin{exercise}[指数率的精确区间]\label{q:B13.1}
指数样本总和为 $T$，容量 $n$。利用 $2\lambda T\sim\chi^2_{2n}$ 构造率 $\lambda$ 与平均寿命 $1/\lambda$ 的 $1-\alpha$ 置信区间。\par\smallskip {\small 详解见第~\pageref{a:B13.1} 页。}\end{exercise}

\begin{exercise}[均匀端点的单侧区间]\label{q:B13.2}
设 $M=X_{(n)}$ 来自 $U(0,\theta)$ 样本。利用 $M/\theta$ 的分布构造形如 $[M,cM]$ 的精确 $1-\alpha$ 置信区间。\par\smallskip {\small 详解见第~\pageref{a:B13.2} 页。}\end{exercise}

\begin{exercise}[两个均值差的合并方差区间]\label{q:B13.3}
两组独立正态样本容量为 $n,m$，均值为 $\mu_1,\mu_2$，共同未知方差为 $\sigma^2$。构造均值差的 $1-\alpha$ 置信区间，并指出共同方差假设的位置。\par\smallskip {\small 详解见第~\pageref{a:B13.3} 页。}\end{exercise}

\begin{exercise}[韦尔奇近似区间]\label{q:B13.4}
两组独立正态样本的方差可能不同。写出均值差的韦尔奇区间及近似自由度。用 $n=10,m=15,s_1^2=4,s_2^2=9,\overline x-\overline y=2$ 给出 $95\%$ 区间的计算式。\par\smallskip {\small 详解见第~\pageref{a:B13.4} 页。}\end{exercise}

\begin{exercise}[配对差的区间]\label{q:B13.5}
对五名受试者作配对测量，所得差值依次为 $2,1,3,0,4$。假定不同受试者相互独立，差值来自正态总体。求平均差的 $95\%$ 置信区间；解释为何不应把前后数据当作两组独立样本。\par\smallskip {\small 详解见第~\pageref{a:B13.5} 页。}\end{exercise}

\begin{exercise}[置信区间与后验区间不能互换]\label{q:M-P10-5}
参数 $\mu$ 只取 $0,1,2,3$，先验概率依次为 $0.94,0.02,0.02,0.02$。条件于 $\mu$，$64$ 个观测独立来自 $N(\mu,25)$，得到 $\overline x=1.6$。构造通常的 $95\%$ 正态均值区间，并计算该已实现区间在指定先验与后验下包含参数的概率。\par\smallskip {\small 详解见第~\pageref{a:M-P10-5} 页。}\end{exercise}

\begin{exercise}[自助法百分位区间]\label{q:B13.6}
给定来自同一总体的独立样本，容量为 $n$。希望估计总体中位数。说明非参数自助法构造百分位区间的步骤、随机性来源以及一种可能失效的情形。\par\smallskip {\small 详解见第~\pageref{a:B13.6} 页。}\end{exercise}

\clearpage\section{假设检验与功效}

\begin{exercise}[硬币的精确双侧检验]\label{q:M-P8-1ab}
独立旋转硬币 $250$ 次，观察到 $140$ 次正面。检验 $H_0:p=0.5$ 对 $H_1:p\ne0.5$。计算精确双侧 $p$ 值，并判断在 $0.10$ 与 $0.05$ 水平是否拒绝。比较不作连续性修正的正态近似。\par\smallskip {\small 详解见第~\pageref{a:M-P8-1ab} 页。}\end{exercise}

\begin{exercise}[离散拒绝域与功效]\label{q:M-P8-1cd}
继续考虑 $X\sim\operatorname{Bin}(250,p)$。构造水平不超过 $0.01$ 的对称双侧拒绝域；再对水平不超过 $0.05$ 的对称拒绝域，分别计算 $p=0.55$ 和 $p=0.60$ 时的功效。\par\smallskip {\small 详解见第~\pageref{a:M-P8-1cd} 页。}\end{exercise}

\begin{exercise}[第一类与第二类错误]\label{q:M-P8-2}
有 $140$ 个真实阴性对象和 $140$ 个真实阳性对象。规则将真实阴性中的 $9$ 个判为阳性，将真实阳性中的 $15$ 个判为阴性。以“对象真实阴性”为零假设，估计第一类错误率、第二类错误率及功效，并解释为什么不能直接得到阳性结果为假的概率。\par\smallskip {\small 详解见第~\pageref{a:M-P8-2} 页。}\end{exercise}

\begin{exercise}[复合零假设与功效规划]\label{q:M-P8-3}
独立仪器读数 $X_i\sim N(\mu,5^2)$，检验 $H_0:\mu\le40$ 对 $H_1:\mu>40$。取 $n=3$，要求所有零假设参数下误拒概率不超过 $0.04$，求样本均值阈值及 $\mu=45$ 时的功效。再求使该功效至少为 $0.9$ 的最小样本量。\par\smallskip {\small 详解见第~\pageref{a:M-P8-3} 页。}\end{exercise}

\begin{exercise}[直接计算两类错误]\label{q:M-P8-4}
只观察一次 $X\sim U(0,\theta)$。检验 $H_0:\theta=2$ 对 $H_1:\theta\ne2$，规定当 $X\le0.1$ 或 $X\ge1.9$ 时拒绝。求第一类错误概率，以及 $\theta=2.5$ 时的第二类错误概率；写出一般功效函数。\par\smallskip {\small 详解见第~\pageref{a:M-P8-4} 页。}\end{exercise}

\begin{exercise}[单侧与双侧 t 检验]\label{q:M-P9-1}
正态样本容量 $16$，样本均值 $11$，无偏样本方差 $4$。检验总体均值为 $10$，显著性水平 $0.05$。分别考虑双侧备择与预先指定的右侧备择，计算统计量、$p$ 值和结论。\par\smallskip {\small 详解见第~\pageref{a:M-P9-1} 页。}\end{exercise}

\begin{exercise}[停止规则改变抽样分布]\label{q:M-P9-2}
观察到 $12$ 次抛掷，其中 $3$ 次正面、$9$ 次反面，最后一次为第三个正面。检验硬币 $p=0.5$ 对 $p<0.5$。分别按“固定抛 $12$ 次”和“直到第三个正面才停止”的计划求单侧 $p$ 值；若采用贝塔先验，比较后验。\par\smallskip {\small 详解见第~\pageref{a:M-P9-2} 页。}\end{exercise}

\begin{exercise}[正态右尾检验的分位数]\label{q:M-E2-II4}
某标准正态右尾检验的水平为 $0.1$，观测统计量为 $z=2$。求临界值、$p$ 值和结论，并说明分布函数与分位数函数分别承担什么作用。\par\smallskip {\small 详解见第~\pageref{a:M-E2-II4} 页。}\end{exercise}

\begin{exercise}[检验与置信区间对偶]\label{q:M-F-C10}
一个双侧两样本均值检验得到 $p=0.08$。假设置信区间与检验使用同一模型和同一统计量，判断：在 $0.05$ 水平拒绝均值相等；$90\%$ 均值差区间包含零；$95\%$ 区间包含零；参数以 $95\%$ 概率属于已计算的 $95\%$ 区间。\par\smallskip {\small 详解见第~\pageref{a:M-F-C10} 页。}\end{exercise}

\begin{exercise}[奈曼—皮尔逊检验]\label{q:B14.1}
独立样本来自 $N(\mu,\sigma^2)$，$\sigma^2$ 已知。检验简单假设 $H_0:\mu=\mu_0$ 对 $H_1:\mu=\mu_1>\mu_0$。利用似然比推导最强水平 $\alpha$ 检验及其功效。\par\smallskip {\small 详解见第~\pageref{a:B14.1} 页。}\end{exercise}

\begin{exercise}[正态均值的似然比检验]\label{q:B14.2}
正态样本方差已知，检验 $H_0:\mu=\mu_0$ 对 $H_1:\mu\ne\mu_0$。求广义似然比 $\Lambda=L(\mu_0)/\sup_\mu L(\mu)$，并说明它如何给出双侧正态检验。\par\smallskip {\small 详解见第~\pageref{a:B14.2} 页。}\end{exercise}

\begin{exercise}[多重检验与邦费罗尼控制]\label{q:B14.3}
同时进行 $20$ 个检验，每个检验在真实零假设下误拒概率为 $0.05$。若全部零假设真实且检验相互独立，求至少一次误拒的概率。若不保证独立，怎样把至少一次误拒的概率控制在 $0.05$？\par\smallskip {\small 详解见第~\pageref{a:B14.3} 页。}\end{exercise}

\begin{exercise}[不拒绝不等于等效]\label{q:B14.4}
某均值差估计为 $0.2$，标准误为 $1$，大样本双侧检验未拒绝差为零。能否据此认定两个总体等效？若预先规定等效界为 $[-0.5,0.5]$，给出一种合适的判断方法。\par\smallskip {\small 详解见第~\pageref{a:B14.4} 页。}\end{exercise}

\clearpage\section{拟合优度与列联表}

\begin{exercise}[单个正态总体的方差检验]\label{q:M-P9-3}
十二个独立正态观测为 \[\begin{gathered}1.76,-2.28,-0.56,1.46,0.59,1.26,\\-1.94,-0.79,-0.86,-1.41,2.07,1.30.\end{gathered}\]。在 $0.05$ 水平检验 $H_0:\sigma^2=1$ 对 $H_1:\sigma^2>1$。\par\smallskip {\small 详解见第~\pageref{a:M-P9-3} 页。}\end{exercise}

\begin{exercise}[两个总体的类别比例比较]\label{q:M-E2-II5}
两组各 $50$ 个独立对象，第一组两类别计数为 $(15,35)$，第二组为 $(5,45)$。利用未作连续性修正的皮尔逊卡方检验，在 $0.05$ 水平检验类别比例是否相同。说明结论的范围。\par\smallskip {\small 详解见第~\pageref{a:M-E2-II5} 页。}\end{exercise}

\begin{exercise}[多项分类的拟合优度]\label{q:B15.1}
公平六面骰子独立掷 $120$ 次，各面计数为 $18,22,19,21,17,23$。用皮尔逊拟合优度检验公平性，写出统计量、自由度和结论。\par\smallskip {\small 详解见第~\pageref{a:B15.1} 页。}\end{exercise}

\begin{exercise}[估计参数后的自由度]\label{q:B15.2}
把来自一个泊松分布的独立计数按预先规定的 $k$ 个互斥类别分组，用样本估计未知率参数，再做拟合优度检验。通常使用什么自由度？说明使用该结论需注意的条件。\par\smallskip {\small 详解见第~\pageref{a:B15.2} 页。}\end{exercise}

\begin{exercise}[列联表独立性检验]\label{q:B15.3}
某 $2\times3$ 表的两行计数为 $(20,30,50)$ 与 $(30,30,40)$。假定样本可按多项模型处理，检验行变量与列变量独立。求期望计数、统计量和自由度。\par\smallskip {\small 详解见第~\pageref{a:B15.3} 页。}\end{exercise}

\begin{exercise}[小样本的费希尔精确检验]\label{q:B15.4}
一个 $2\times2$ 表的第一行为 $(4,1)$，第二行为 $(1,4)$。在固定行列边际和的条件下，以第一行第一格 $A$ 为统计量，求右侧精确 $p$ 值和按概率不大于观察表概率定义的双侧 $p$ 值。\par\smallskip {\small 详解见第~\pageref{a:B15.4} 页。}\end{exercise}

\begin{exercise}[两个正态方差的 F 检验]\label{q:M-P9-5}
两组独立正态样本各有 $20$ 个观测，样本方差分别为 $1.1302$ 与 $1.1647$。检验 $H_0:\sigma_1^2=\sigma_2^2$ 对双侧备择，计算 $F$ 统计量、双侧 $p$ 值并构造方差比的 $95\%$ 区间。\par\smallskip {\small 详解见第~\pageref{a:M-P9-5} 页。}\end{exercise}

\clearpage\section{回归与方差分析}

\begin{exercise}[三个点的最小二乘直线]\label{q:M-C26-1a}
给定三个点 $(1,3)$、$(2,1)$、$(4,4)$。写出含截距的直线模型，最小化残差平方和求拟合直线，并验证拟合线经过样本均值点。\par\smallskip {\small 详解见第~\pageref{a:M-C26-1a} 页。}\end{exercise}

\begin{exercise}[二次多项式仍是线性回归]\label{q:M-C26-1b}
对点 $(1,3)$、$(2,1)$、$(4,4)$ 拟合 $y=ax^2+bx+c$。给出设计矩阵和解，并解释为什么二次、三次曲线也属于线性回归。三个观测拟合三个参数有什么统计限制？\par\smallskip {\small 详解见第~\pageref{a:M-C26-1b} 页。}\end{exercise}

\begin{exercise}[对数尺度的指数拟合]\label{q:M-C26-1c}
仍使用三个正响应点 $(1,3)$、$(2,1)$、$(4,4)$，在对数尺度拟合 $\log y=ax+b$。写出最小化目标和参数公式，并解释它与直接最小化原尺度误差 $\sum_i(y_i-e^{ax_i+b})^2$ 是否相同。\par\smallskip {\small 详解见第~\pageref{a:M-C26-1c} 页。}\end{exercise}

\begin{exercise}[正态误差下的最小二乘与似然]\label{q:M-C26-2d}
在固定设计模型 $Y_i=ax_i+b+\varepsilon_i$ 中，误差独立同分布于 $N(0,\sigma^2)$。证明 $a,b$ 的极大似然估计等于最小二乘估计，并求误差方差的极大似然估计。\par\smallskip {\small 详解见第~\pageref{a:M-C26-2d} 页。}\end{exercise}

\begin{exercise}[正规方程与残差正交]\label{q:B16.1}
线性模型写成 $Y=X\beta+\varepsilon$，其中固定设计矩阵 $X$ 为 $n\times p$ 满列秩矩阵。推导最小二乘解，证明残差与设计矩阵各列正交，并说明包含截距时残差和为零。\par\smallskip {\small 详解见第~\pageref{a:B16.1} 页。}\end{exercise}

\begin{exercise}[高斯—马尔可夫最优性]\label{q:B16.2}
在固定满列秩设计模型中，假设 $E\varepsilon=0$、$\operatorname{Cov}(\varepsilon)=\sigma^2I$。证明最小二乘估计在线性无偏估计中协方差最小，不要求误差正态。\par\smallskip {\small 详解见第~\pageref{a:B16.2} 页。}\end{exercise}

\begin{exercise}[斜率估计的方差与区间]\label{q:B16.3}
在含截距的一元正态线性回归中，记 $S_{xx}=\sum_i(x_i-\overline x)^2>0$。推导斜率估计的均值、方差及置信区间，并说明扩大自变量跨度的作用。\par\smallskip {\small 详解见第~\pageref{a:B16.3} 页。}\end{exercise}

\begin{exercise}[均值响应区间与预测区间]\label{q:B16.4}
在一元正态线性回归中，对固定 $x_0$，分别构造均值响应 $a x_0+b$ 的区间和一个独立新响应 $Y_0$ 的预测区间。解释额外的方差项。\par\smallskip {\small 详解见第~\pageref{a:B16.4} 页。}\end{exercise}

\begin{exercise}[单因素方差分析]\label{q:B16.5}
三组独立正态观测分别为 $(2,3,4)$、$(4,5,6)$、$(6,7,8)$，假定共同方差。用单因素方差分析在 $0.05$ 水平检验三个总体均值相等。\par\smallskip {\small 详解见第~\pageref{a:B16.5} 页。}\end{exercise}

\begin{exercise}[相关系数与回归解释度]\label{q:B16.6}
对含截距的一元最小二乘回归，证明拟合平方和与残差平方和正交分解，并证明 $R^2=r^2$。说明该等式需要排除哪些退化情形。\par\smallskip {\small 详解见第~\pageref{a:B16.6} 页。}\end{exercise}

\clearpage\section{综合训练与解题辨析}

\begin{exercise}[帕累托模型从估计到区间]\label{q:M-F-5}
独立观测来自密度 $f(x;\theta)=\theta x^{-\theta-1}$，$x\ge1,\theta>0$。（1）求极大似然估计；（2）证明 $Y=\log X$ 服从指数分布；（3）构造 $\theta$ 的精确置信区间；（4）说明总体均值何时存在。\par\smallskip {\small 详解见第~\pageref{a:M-F-5} 页。}\end{exercise}

\begin{exercise}[瑞利参数与变换统计量]\label{q:M-PF-26}
样本密度为 $f(x;\tau)=\tau x\exp(-\tau x^2/2)$，$x>0,\tau>0$。（1）求 $\tau$ 的极大似然估计；（2）求 $Z=X^2/2$ 的分布；（3）给出无偏率估计与精确置信区间。\par\smallskip {\small 详解见第~\pageref{a:M-PF-26} 页。}\end{exercise}

\begin{exercise}[定位中的连续更新与预测]\label{q:M-F-7}
四扇门依次为 $A,B,C,D$，颜色分别为橙、橙、蓝、蓝。机器人初始等概率位于各门；扫描颜色正确率为 $0.7$。首次读到橙色后，求位置后验；随后确定地向右移一扇门（$D$ 后接 $A$），求新位置先验和下一次读到蓝色的预测概率。各次扫描在给定真实位置后独立。\par\smallskip {\small 详解见第~\pageref{a:M-F-7} 页。}\end{exercise}

\begin{exercise}[六类别概率的综合检验]\label{q:M-F-8}
一个教学模型声称六类概率为 $(0.20,0.20,0.20,0.15,0.15,0.10)$。独立抽取 $1000$ 个对象，观察计数为 $(180,190,185,\allowbreak160,165,120)$。在 $0.05$ 水平检验模型，求统计量、自由度与结论，并指出若对象按包装聚集抽样可能出现的问题。\par\smallskip {\small 详解见第~\pageref{a:M-F-8} 页。}\end{exercise}

\begin{exercise}[两样本检验的完整流程]\label{q:M-F-9}
甲乙两组各有 $11$ 个独立观测，样本均值分别为 $115,110$，样本标准差为 $8,6$。在正态且共同方差的模型下进行双侧均值检验，显著性水平 $0.05$；再求甲总体均值的 $95\%$ 置信区间。\par\smallskip {\small 详解见第~\pageref{a:M-F-9} 页。}\end{exercise}

\begin{exercise}[双侧检验在指定备择下的功效]\label{q:M-F-11}
独立抽取 $100$ 个 $N(\mu,1)$ 观测，以 $0.05$ 水平检验 $H_0:\mu=0$ 对双侧备择。求当 $\mu=0.4$ 时的功效，并说明原资料将接受区间写成拒绝区间会造成什么错误。\par\smallskip {\small 详解见第~\pageref{a:M-F-11} 页。}\end{exercise}

\begin{exercise}[折线基函数的最小二乘]\label{q:M-F-12}
观测为 $(4,1)$、$(-2,1.5)$、$(0,0.5)$，模型为 $Y_i=b|x_i-2|+\varepsilon_i$。求最小二乘估计；在独立零均值同方差误差下求其方差；描述拟合曲线的形状。\par\smallskip {\small 详解见第~\pageref{a:M-F-12} 页。}\end{exercise}

\begin{exercise}[自助分位数的读法]\label{q:M-F-13}
原样本容量 $36$，均值 $7.0$，标准差 $0.84$。进行 $500$ 次有放回自助重抽样并计算样本均值。排序后第 $25$ 个值为 $6.740$，第 $475$ 个值为 $7.224$。采用经验分布逆函数的分位数约定，构造 $90\%$ 百分位自助区间，并讨论重抽样次数的影响。\par\smallskip {\small 详解见第~\pageref{a:M-F-13} 页。}\end{exercise}

\begin{exercise}[抽球机制与逆向条件概率]\label{q:M-PF-18}
袋中有 $3$ 个橙球、$2$ 个蓝球。第一次若取到橙球，则取出不放回；若取到蓝球，则放回并另加一个橙球。随后取第二个球。求两次均橙的概率；已知第二次为橙，求第一次为蓝的概率。\par\smallskip {\small 详解见第~\pageref{a:M-PF-18} 页。}\end{exercise}

\begin{exercise}[停时分布的边界计算]\label{q:M-PF-19}
反复独立投公平六面骰子，直到累计点数严格大于 $6$ 才停止，记投掷次数为 $N$。求 $P(N\le1),P(N\le2),P(N\le7)$，再利用尾和求 $EN$。\par\smallskip {\small 详解见第~\pageref{a:M-PF-19} 页。}\end{exercise}

\begin{exercise}[分段密度与分位数]\label{q:M-PF-20}
成绩 $X$ 在 $[0,1]$ 上具有密度：$0\le x\le1/2$ 时为 $4x$，$1/2<x\le1$ 时为 $4-4x$。求及格线为 $0.55$ 时的通过概率、$87.5\%$ 分位数，并写出整个分布函数。\par\smallskip {\small 详解见第~\pageref{a:M-PF-20} 页。}\end{exercise}

\begin{exercise}[从分布函数回到密度与期望]\label{q:M-PF-21}
分布函数在 $x<0$ 时为零，在 $0\le x\le1$ 时为 $x(2-x)$，在 $x>1$ 时为一。求密度、$EX$、$\operatorname{Var}(X)$ 和 $P(X<0.4)$。说明端点是否有原子概率。\par\smallskip {\small 详解见第~\pageref{a:M-PF-21} 页。}\end{exercise}

\begin{exercise}[参数自助法与基本区间]\label{q:M-PF-29}
连续独立伯努利试验的成功概率为 $p$，记录直到首次失败的总试验次数 $X$（包括失败），故 $P(X=k)=p^{k-1}(1-p)$。给定容量 $n$ 的样本，估计 $p$，并说明参数自助法构造 $95\%$ 基本区间的步骤。\par\smallskip {\small 详解见第~\pageref{a:M-PF-29} 页。}\end{exercise}

\begin{exercise}[极小尾概率与数值稳定性]\label{q:M-PF-30}
独立抽取 $100$ 个 $N(\mu,4)$ 观测，双侧检验 $H_0:\mu=3$，水平 $0.05$。若样本均值为 $5$，求 $p$ 值；再求 $\mu=4$ 时的功效。解释为什么不能把很小的正态尾概率写成严格的零。\par\smallskip {\small 详解见第~\pageref{a:M-PF-30} 页。}\end{exercise}

\begin{exercise}[正态测量的后验精度规划]\label{q:M-PF-31}
未知量 $\theta$ 的先验为 $N(80,4)$，每次测量误差独立且服从 $N(0,0.5)$。一次读数为 $85$，求后验分布。若继续独立测量，总共至少测量多少次可使后验方差小于 $0.01$？\par\smallskip {\small 详解见第~\pageref{a:M-PF-31} 页。}\end{exercise}

\begin{exercise}[泊松分裂与参数推断]\label{q:B17.1}
一天到达的对象数 $N\sim\operatorname{Pois}(\lambda)$，每个对象独立以概率 $p$ 标为甲类。记两类数量为 $X,Y$。（1）推导联合分布并判断独立性；（2）连续观察 $m$ 天两类总数 $T_X,T_Y$，求 $\lambda,p$ 的极大似然估计；（3）说明零总数时的问题。\par\smallskip {\small 详解见第~\pageref{a:B17.1} 页。}\end{exercise}

\begin{exercise}[有限总体抽样与修正因子]\label{q:B17.2}
有限总体有 $N$ 个对象，其中 $M$ 个具有某特征。无放回简单随机抽样 $n$ 个，令 $X$ 为特征数、$\widehat p=X/n$。求分布、均值和方差，比较与有放回伯努利抽样的差别。\par\smallskip {\small 详解见第~\pageref{a:B17.2} 页。}\end{exercise}

\begin{exercise}[随机化回答的估计误差]\label{q:B17.3}
为保护隐私，每名对象独立掷一枚公平硬币：正面时回答真实的是非问题，反面时强制回答“是”。真实“是”的总体比例为 $p$。对 $n$ 人记录“是”的比例 $\widehat q$。构造 $p$ 的无偏估计，求方差并解释参数越界问题。\par\smallskip {\small 详解见第~\pageref{a:B17.3} 页。}\end{exercise}

\begin{exercise}[尾和公式与收集全套]\label{q:B17.4}
独立抽取 $m$ 种等概率卡片，每次一张。记收集到全部 $m$ 种所需次数为 $T$。把过程分解为等待新类型的阶段，求 $ET$ 和 $\operatorname{Var}(T)$，并说明独立性依据。\par\smallskip {\small 详解见第~\pageref{a:B17.4} 页。}\end{exercise}

\begin{exercise}[随机和的复合泊松结构]\label{q:B17.5}
设 $N\sim\operatorname{Pois}(\lambda)$，$Z_i$ 独立同分布且与 $N$ 独立，$EZ_i=\mu$、$\operatorname{Var}(Z_i)=\sigma^2$。令 $S=\sum_{i=1}^NZ_i$，空和为零。求 $ES,\operatorname{Var}(S)$ 和特征函数，并说明只给前两矩是否足以精确求 $P(S\le c)$。\par\smallskip {\small 详解见第~\pageref{a:B17.5} 页。}\end{exercise}

\begin{exercise}[正态变量平方和与比值]\label{q:B17.6}
设 $Z_1,Z_2,Z_3$ 独立标准正态。令 $U=Z_1^2+Z_2^2$、$V=Z_3^2$。求 $U,V$ 的分布、$U/(2V)$ 的分布，以及 $U/(U+V)$ 的分布。说明后两个比值的不同支撑。\par\smallskip {\small 详解见第~\pageref{a:B17.6} 页。}\end{exercise}

\begin{exercise}[指数竞争与获胜标记]\label{q:B17.7}
设 $X\sim\operatorname{Exp}(\lambda)$、$Y\sim\operatorname{Exp}(\mu)$ 独立。令 $T=\min(X,Y)$、$J=\mathbf1_{\{X<Y\}}$。求 $T,J$ 的分布并证明独立性。\par\smallskip {\small 详解见第~\pageref{a:B17.7} 页。}\end{exercise}

\begin{exercise}[条件期望作为最小平方预测]\label{q:B17.8}
设 $EY^2<\infty$，以 $X$ 的可测函数 $g(X)$ 预测 $Y$。证明 $E(Y\mid X)$ 最小化均方预测误差，并给出误差分解。若 $(X,Y)$ 联合正态，写出最优预测。\par\smallskip {\small 详解见第~\pageref{a:B17.8} 页。}\end{exercise}

\begin{exercise}[混合正态的方差分解]\label{q:B17.9}
先以概率 $w$ 选第一组、概率 $1-w$ 选第二组，随后分别从 $N(\mu_1,\sigma_1^2)$ 或 $N(\mu_2,\sigma_2^2)$ 生成 $X$。求均值和方差。若两个组的均值相同，混合分布是否必为正态？\par\smallskip {\small 详解见第~\pageref{a:B17.9} 页。}\end{exercise}

\begin{exercise}[德尔塔方法估计风险函数]\label{q:B17.10}
独立伯努利样本参数为 $p\in(0,1)$，样本比例为 $\widehat p$。推导对数赔率 $g(p)=\log[p/(1-p)]$ 的代入估计渐近分布，构造近似区间，并解释边界观测问题。\par\smallskip {\small 详解见第~\pageref{a:B17.10} 页。}\end{exercise}

\begin{exercise}[威尔逊比例区间的反演]\label{q:B17.11}
对伯努利比例 $p$，从正态得分不等式 $|\widehat p-p|\le z\sqrt{p(1-p)/n}$ 出发推导威尔逊区间。指出它与把 $p$ 直接替成 $\widehat p$ 的沃尔德区间有何区别。\par\smallskip {\small 详解见第~\pageref{a:B17.11} 页。}\end{exercise}

\begin{exercise}[正态模型的充分统计量]\label{q:B17.12}
正态样本的 $\mu,\sigma^2$ 都未知。证明 $(\sum_iX_i,\sum_iX_i^2)$ 是充分统计量，说明它与 $(\overline X,S^2)$ 的等价关系，并讨论只保留均值是否足够。\par\smallskip {\small 详解见第~\pageref{a:B17.12} 页。}\end{exercise}

\begin{exercise}[最小值模型的非对称估计]\label{q:B17.13}
设 $X_i=\theta+E_i$，其中 $E_i$ 独立服从率 $\lambda$ 已知的指数分布，$\theta\in\mathbb R$。求 $\theta$ 的极大似然估计、无偏修正和一个精确 $1-\alpha$ 置信区间。\par\smallskip {\small 详解见第~\pageref{a:B17.13} 页。}\end{exercise}

\begin{exercise}[贝塔—二项预测的额外变异]\label{q:B17.14}
给定参数 $p$ 后，未来 $m$ 次试验独立伯努利；当前后验为 $p\mid D\sim\operatorname{Beta}(a,b)$。令未来成功数为 $K$。求预测质量函数、均值和方差，并与固定 $p=a/(a+b)$ 的二项分布比较。\par\smallskip {\small 详解见第~\pageref{a:B17.14} 页。}\end{exercise}

\begin{exercise}[正态性、独立性与零相关]\label{q:B17.15}
令 $Z\sim N(0,1)$，独立符号变量 $S$ 等概率取 $-1,1$。定义 $X=Z,Y=SZ$。证明两个边缘都为标准正态且协方差为零，但不独立。说明这为何不违反联合正态中的零相关结论。\par\smallskip {\small 详解见第~\pageref{a:B17.15} 页。}\end{exercise}

\begin{exercise}[支撑集破坏独立性]\label{q:B17.16}
设 $(X,Y)$ 在三角形 $0<x<y<1$ 上密度为 $2$，其余处为零。有人因内部密度可写成常数乘常数而认定独立。计算边缘密度、协方差并指出错误。\par\smallskip {\small 详解见第~\pageref{a:B17.16} 页。}\end{exercise}

\begin{exercise}[蒙特卡洛积分与标准误]\label{q:B17.17}
用独立 $U_i\sim U(0,1)$ 估计积分 $I=\int_0^1x^2dx$，取 $\widehat I=n^{-1}\sum_iU_i^2$。求偏差、方差、近似置信区间与切比雪夫误差界；比较误差随计算量的变化。\par\smallskip {\small 详解见第~\pageref{a:B17.17} 页。}\end{exercise}

\begin{exercise}[反向变量的方差缩减]\label{q:B17.18}
仍估计 $\int_0^1x^2dx$。对一个均匀变量 $U$，用 $A=[U^2+(1-U)^2]/2$ 作为一次配对估计。求其期望和方差，与使用两个独立均匀变量的平均比较。\par\smallskip {\small 详解见第~\pageref{a:B17.18} 页。}\end{exercise}

\begin{exercise}[切比雪夫界与正态近似的角色]\label{q:B17.19}
独立同分布样本均值为 $\mu$，方差为 $4$，容量 $100$。估计 $P(|\overline X-\mu|<0.5)$。分别给出仅用方差的下界、中心极限定理近似，以及总体正态时的精确值。\par\smallskip {\small 详解见第~\pageref{a:B17.19} 页。}\end{exercise}

\begin{exercise}[学生化与斯卢茨基定理]\label{q:B17.20}
设独立同分布样本具有有限正方差 $\sigma^2$，但不要求正态。证明 $\sqrt n(\overline X-\mu)/S\xrightarrow dN(0,1)$，并解释它与有限样本 $t_{n-1}$ 分布的区别。\par\smallskip {\small 详解见第~\pageref{a:B17.20} 页。}\end{exercise}

\begin{exercise}[三种收敛之间的辨析]\label{q:B17.21}
证明均方收敛蕴含依概率收敛。给出一个依概率收敛到零但期望不趋于零的非负随机变量序列，并解释这对交换极限与期望的警示。\par\smallskip {\small 详解见第~\pageref{a:B17.21} 页。}\end{exercise}

\begin{exercise}[参数空间约束下的正态估计]\label{q:B17.22}
样本来自 $N(\mu,1)$，已知 $\mu\ge0$。求约束极大似然估计，并比较它与无约束均值估计在平方损失下的误差。它是否仍然无偏？\par\smallskip {\small 详解见第~\pageref{a:B17.22} 页。}\end{exercise}

\begin{exercise}[遗漏变量导致斜率变化]\label{q:B17.23}
总体满足 $Y=\beta_0+\beta_1X+\beta_2Z+\varepsilon$，其中 $E\varepsilon=0$，误差与 $X$ 不相关，二阶矩有限且 $\operatorname{Var}(X)>0$。如果只用 $X$ 对 $Y$ 做含截距的总体最小平方线性预测，斜率是多少？何时等于 $\beta_1$？\par\smallskip {\small 详解见第~\pageref{a:B17.23} 页。}\end{exercise}

\begin{exercise}[加权最小二乘]\label{q:B17.24}
模型为 $Y=X\beta+\varepsilon$，$X$ 满列秩，$E\varepsilon=0$，协方差为已知正定矩阵 $\Sigma$。推导广义最小二乘估计及其协方差。若误差独立且方差不同，解释各观测的权重。\par\smallskip {\small 详解见第~\pageref{a:B17.24} 页。}\end{exercise}

\clearpage